\chapter{Алгебра}\label{chap:algebra}

В предыдущей главе мы дали введение в базовые вычислительные инструменты, необходимые для подхода к SNARKам «бумага и ручка». В этой главе мы дадим более абстрактное пояснение соответствующей математической терминологии, такой как \term{группы} (groups), \term{кольца} (rings) и \term{поля} (fields).

Научная литература по криптографии часто содержит такие термины, и необходимо получить хотя бы общее представление о них, чтобы иметь возможность следить за изложением.

\section{Коммутативные группы}
\label{sec:groups}
Коммутативные группы — это абстракции, которые отражают суть математических явлений, таких как сложение и вычитание, или умножение и деление.

Чтобы понять коммутативные группы, вспомним, как в школе мы изучали сложение и вычитание целых чисел. Мы узнали, что всякий раз, когда складываем два целых числа, результат гарантированно тоже целое. Мы также узнали, что прибавление нуля к любому целому числу означает, что «ничего не происходит», так как результат сложения — это то же самое целое число. Кроме того, мы узнали, что порядок, в котором мы складываем два (или более) целых числа, не имеет значения, что скобки не влияют на результат сложения, и что для каждого целого числа всегда есть другое целое число (отрицательное), такое что при их сложении получается ноль.

Эти условия являются определяющими свойствами коммутативной группы, и математики осознали, что точно такой же набор правил можно найти в совершенно разных математических структурах. Поэтому имеет смысл дать абстрактное, формальное определение того, что такое группа, отвязавшись от конкретных примеров, таких как целые числа. Это позволяет нам гибко работать с объектами самого разного математического происхождения, сохраняя при этом существенные структурные аспекты многих объектов абстрактной алгебры и не только.

Выделяя эти правила в минимальный независимый список свойств и делая их абстрактными, мы приходим к следующему определению коммутативной группы:

\tbds{revise the counter for definitions, current one too long}

\begin{definition}
\label{def:commutative_group}
\term{Коммутативная группа} $(\G,\cdot) $ состоит из множества $\G$ и \term{отображения} $\cdot:\G \times \G \to \G $. Отображение называется \term{групповым законом} (group law), и оно комбинирует два элемента множества $ \G$ в третий так, что выполняются следующие свойства:
\begin{itemize}
\item \hilight{Коммутативность} (Commutativity): Для всех $g_1,g_2\in\G$ выполняется равенство $g_1\cdot g_2=g_2\cdot g_1$.
\item \hilight{Ассоциативность} (Associativity): Для всех $g_1,g_2,g_3\in\G$ выполняется равенство
$g_1\cdot(g_2\cdot g_3) = (g_1\cdot g_2)\cdot g_3$.
\item \hilight{Существование нейтрального элемента} (Existence of a neutral element): Для каждого $g\in\G$ существует $e\in\G$ такой, что $e\cdot g=g$.
\item \hilight{Существование обратного элемента} (Existence of an inverse): Для каждого $g\in\G$ существует $g^{-1}\in\G$ такой, что $g\cdot g^{-1}=e$.
\end{itemize}
Если $(\G,\cdot)$ — группа, и $\G'\subseteq\G$ — подмножество $\G$ такое, что \term{ограничение} группового закона $\cdot: \G'\times \G' \to \G'$ является групповым законом на $\G'$, то $(\G',\cdot)$ называется \term{подгруппой} (subgroup) группы $(\G,\cdot)$.
\end{definition}

Перефразируя абстрактное определение простым языком, группа — это нечто, где мы можем выполнять вычисления способом, напоминающим поведение сложения целых чисел. В частности, это означает, что мы можем комбинировать некоторый элемент с другим элементом в новый элемент так, что эта операция обратима и порядок комбинирования элементов не имеет значения.

\begin{notation}Поскольку в этой книге мы рассматриваем исключительно коммутативные группы, мы часто просто называем их группами, оставляя коммутативность неявной.\footnote{Коммутативные группы также называются \term{абелевыми группами} (Abelian groups). Множество $\G$ с отображением $\cdot$, удовлетворяющим всем перечисленным выше правилам, кроме закона коммутативности, называется \term{некоммутативной группой} (non-commutative group).} 

Если нет риска неоднозначности (относительно того, каков групповой закон данной группы), мы часто опускаем символ $\cdot$ и просто пишем $\mathbb{G}$ как обозначение группы, оставляя групповой закон неявным. В этом случае мы также говорим, что $\G$ имеет групповой тип, указывая на то, что $\G$ — это не просто множество, а множество вместе с групповым законом.
\end{notation}

\begin{notation}[\deftitle{Аддитивная запись}]\label{def:additive_notation}
Для коммутативных групп $(\G,\cdot)$ мы иногда используем так называемую \term{аддитивную запись} (additive notation) $(\G,+)$, то есть пишем $+$ вместо $\cdot$ для группового закона, $0$ для нейтрального элемента и $-g:=g^{-1}$ для обратного элемента $g\in\G$.
\end{notation}
Как мы увидим в последующих главах, группы широко используются в криптографии и в SNARKах.\footnote{Более углублённое введение в коммутативные группы можно найти, например, в \chaptname{} 1, \secname{} 1 книги \cite{nieder-1986} или в \chaptname{} 1 книги \cite{fuchs-2015}. Введение, более ориентированное на нужды криптографии, можно найти, например, в \chaptname{} 3, \secname{} 8.1.3 книги \cite{katz-2007}.} Но сначала рассмотрим несколько более знакомых примеров.

\begin{example}[Сложение и вычитание целых чисел]
\label{example:group_of_integers}
Множество $(\Z,+)$ целых чисел со сложением — это архетипический пример коммутативной группы, где групповой закон традиционно записывается в аддитивной нотации (\notationname{} \ref{def:additive_notation}). 

Чтобы сравнить сложение целых чисел с абстрактными аксиомами коммутативной группы, отметим сначала, что сложение целых чисел \hilight{коммутативно и ассоциативно}, поскольку $a+b = b+a$, а также $(a+b)+c=a+(b+c)$ для всех целых чисел $a,b,c\in \Z$. \hilight{Нейтральным элементом} $e$ является число $0$, так как $a+0=a$ для всех целых чисел $a\in \Z$. Кроме того, \hilight{обратным} элементом для числа является его отрицательное значение, так как $a+(-a)=0$ для всех $a\in\Z$. Это означает, что целые числа со сложением действительно являются коммутативной группой в абстрактном смысле.

В качестве примера подгруппы группы целых чисел рассмотрим множество чётных чисел, включая $0$. 

$$\Z_{even}:=\{\ldots,-4,-2,0,2,4,\ldots\}$$

Мы видим, что это множество является подгруппой $(\Z,+)$, поскольку сумма двух чётных чисел всегда снова чётна, нейтральный элемент $0$ принадлежит $\Z_{even}$, и отрицательное значение чётного числа само является чётным числом. 
\end{example}
\begin{example}[Тривиальная группа]\label{example:trivial_group}
Самый простой пример коммутативной группы — группа с единственным элементом $\{\bullet\}$ и групповым законом $\bullet\cdot \bullet=\bullet$. Мы называем её \term{тривиальной группой} (trivial group).

Тривиальная группа является подгруппой любой группы. Чтобы увидеть это, пусть $(\G,\cdot)$ — группа с нейтральным элементом $e\in\G$. Тогда $e\cdot e = e$, а также $e^{-1}=e$. Следовательно, множество $\{e\}$ является подгруппой $\G$. В частности, $\{0\}$ является подгруппой $(\Z,+)$, так как $0+0=0$.
\end{example}

\begin{example} Рассмотрим сложение в арифметике по модулю $6$ $(\Z_6,+)$, как определено в примере \ref{def_residue_ring_z_6}. Как мы видим, остаток $0$ является нейтральным элементом в сложении по модулю $6$, а обратным элементом для остатка $r$ является $6-r$, потому что $r+(6-r)=6$. $6$ конгруэнтно $0$, так как $\Zmod{6}{6}=0$. Более того, $r_1+r_2 = r_2 + r_1$, а также $(r_1+r_2)+r_3=r_1+(r_2+r_3)$ наследуются от сложения целых чисел. Поэтому мы видим, что $(\Z_6,+)$ — группа.
\end{example}
Предыдущий пример коммутативной группы очень важен для этой книги. Абстрагируясь от этого примера и рассматривая классы вычетов $(\Z_n,+)$ для произвольных модулей $n$, можно показать, что $(\Z_n,+)$ является коммутативной группой с нейтральным элементом $0$ и аддитивным обратным $n-r$ для любого элемента $r\in\Z_n$. Мы называем такую группу \term{группой классов остатков} (remainder class group) модуля $n$.
\begin{exercise}\label{fstar} Рассмотрим снова пример \ref{primfield_z_5}, и пусть $\Z_5^*$ — множество всех классов остатков из $\Z_5$ без класса $0$. Тогда $\Z_5^*=\{1,2,3,4\}$. Покажите, что $(\Z_5^*,\cdot)$ является коммутативной группой.
\end{exercise}
\begin{exercise}\label{ex:Zn*} Обобщая предыдущее упражнение, рассмотрим общий модуль $n$, и пусть $\Z_n^*$ — множество всех классов остатков из $\Z_n$ без класса $0$. Тогда $\Z_n^*=\{1,2,\ldots,n-1\}$. Приведите контрпример, показывающий, что $(\Z^*_n,\cdot)$ не является группой вообще.

Найдите условие, при котором $(\Z^*_n,\cdot)$ является коммутативной группой, вычислите нейтральный элемент, дайте замкнутую форму для обратного элемента любого элемента и докажите аксиомы коммутативной группы.
\end{exercise}
\subsection{Конечные группы}
\label{sec:finite-groups}
 Как мы видели в предыдущих примерах, группы могут содержать либо бесконечно много элементов (как целые числа), либо конечное количество элементов (как, например, группы классов остатков $(\Z_n,+)$). Чтобы зафиксировать это различие, группа называется \term{конечной группой} (finite group), если лежащее в её основе множество элементов конечно. В этом случае число элементов этой группы называется её \term{порядком} (order).\footnote{Введение в конечные группы дано в \chaptname{} 1 книги \cite{fuchs-2015}. Введение с точки зрения криптографии можно найти в \chaptname{} 3, \secname{} 8.3.1 книги \cite{katz-2007}.}
\begin{notation}
Пусть $\mathbb{G}$ — конечная группа. Мы пишем $ord(\mathbb{G})$ или $|\mathbb{G}|$ для порядка $\mathbb{G}$.
\end{notation}
\begin{example}\label{example:Zn}
Рассмотрим группы классов остатков $(\Z_6,+)$ из примера \ref{def_residue_ring_z_6}, группу $(\Z_5,+)$ из примера \ref{primfield_z_5}, и группу $(\Z_5^*,\cdot)$ из упражнения \ref{fstar}. Мы легко видим, что порядок $(\Z_6,+)$ равен $6$, порядок $(\Z_5,+)$ равен $5$ и порядок $(\Z_5^*,\cdot)$ равен $4$.
\end{example}
\begin{exercise}
\label{ex:Zn} Пусть $n\in\N$ при $n\geq 2$ — некоторый модуль. Каков порядок группы классов остатков $(\Z_n,+)$?
\end{exercise}
\subsection{Образующие}\label{generators} Перечисление множества элементов группы может быть {сложным}, и не всегда очевидно, как фактически вычислить элементы данной группы. С практической точки зрения поэтому желательно иметь группы с \term{множеством образующих} (generator set). Это небольшое подмножество элементов, из которых все остальные элементы могут быть получены применением группового закона многократно только к элементам множества образующих и/или их обратным элементам. 

Конечно, у каждой группы $\G$ есть тривиальное множество образующих, если мы просто считаем каждый элемент группы входящим в множество образующих. Более интересный вопрос — найти наименьшее возможное множество образующих для данной группы. Особого интереса в этом отношении заслуживают группы, у которых множество образующих содержит только один элемент. В этом случае существует (не обязательно единственный) элемент $g\in\G$ такой, что любой другой элемент из $\G$ может быть вычислен повторным комбинированием только $g$ и его обратного элемента $g^{-1}$. 

\begin{definition}[\deftitle{Циклические группы}]\label{cyclic-groups}
Группа $\G$ называется \term{циклической группой} (cyclic group), если существует по крайней мере один элемент $g\in \G$, называемый \term{образующим} (generator), который может породить каждый элемент $G$.
\end{definition}

\begin{example}
\label{example:cyclic_group_of_integers} Самый простой пример циклической группы — группа целых чисел со сложением $(\Z,+)$. В этом случае число $1$ является образующим $\Z$, поскольку каждое целое число может быть получено повторным прибавлением к самому себе либо $1$, либо его обратного элемента $-1$. Например, $-4$ порождается $1$, так как $-4=-1+(-1)+(-1)+(-1)$. Другим образующим $\Z$ является число $-1$.
\end{example}
\begin{example}
\label{example:cyclic_group_F5*} Рассмотрим группу $(\Z_5^*,\cdot)$ из упражнения \ref{fstar}. Поскольку $2^1=2$, $2^2=4$, $2^3=3$ и $2^4=1$, элемент $2$ является образующим $(\Z_5^*,\cdot)$. Кроме того, поскольку $3^1=3$, $3^2=4$, $3^3=2$ и $3^4=1$, элемент $3$ является другим образующим $(\Z_5^*,\cdot)$. Таким образом, циклические группы могут иметь более одного образующего. Однако поскольку $4^1=4$, $4^2=1$, $4^3=4$ и вообще $4^k=4$ для нечётных $k$ и $4^k=1$ для чётных $k$, элемент $4$ не является образующим $(\Z_5^*,\cdot)$. Следовательно, вообще говоря, не каждый элемент конечной циклической группы является образующим.
\end{example}
\begin{example} Рассмотрим модуль $n$ и группы классов остатков $(\Z_n,+)$ из упражнения \ref{ex:Zn}. Эти группы цикличны с образующим $1$, поскольку каждый другой элемент этой группы может быть построен повторным прибавлением к самому себе класса остатков $1$. Поскольку $\Z_n$ также конечна, мы знаем, что $(\Z_n,+)$ — конечная циклическая группа порядка $n$.
\end{example}
\begin{exercise}
\label{example:cyclic_group_F6}
Рассмотрим группу $(\Z_6,+)$ сложения по модулю $6$ из примера \ref{def_residue_ring_z_6}. Покажите, что $5\in \Z_6$ является образующим, а затем покажите, что $2\in \Z_6$ не является образующим.
\end{exercise}
\begin{exercise}\label{ex:modulus-prime-group}(Сложно) Пусть $p\in\mathbb{P}$ — простое число и $(\Z_p^*,\cdot)$ — конечная группа из упражнения \ref{ex:Zn*}. Покажите, что $(\Z_p^*,\cdot)$ циклична.
\end{exercise}

\subsection{Экспоненциальные отображения}
Заметим, что когда $\G$ — циклическая группа порядка $n$ и $g\in \G$ — образующий $\G$, то существует так называемое \textbf{экспоненциальное отображение}, которое отображает аддитивный групповой закон группы классов остатков $(\Z_n,+)$ на групповой закон $\G$ во взаимно однозначном соответствии. Экспоненциальное отображение можно формализовать, как в \eqref{exponentialmap} ниже (где $g^x$ означает «умножить $g$ на самого себя $x$ раз» и $g^0=e_{\G}$).

\begin{equation}\label{exponentialmap}
g^{(\cdot)}: \Z_n \to \G\;;\; x \mapsto g^x
\end{equation}

Чтобы понять, как работает экспоненциальное отображение, сначала заметим, что, поскольку $g^0:=e_{\G}$ по определению, нейтральный элемент $\Z_n$ отображается в нейтральный элемент $\G$. Кроме того, поскольку $g^{x+y}=g^x\cdot g^y$, отображение уважает групповой закон.

\begin{notation}[\deftitle{Скалярное умножение}]
\label{def:scalar_multiplication} Если группа $(\G,+)$ записана в аддитивной нотации (\notationname{} \ref{def:additive_notation}), то экспоненциальное отображение часто называют \term{скалярным умножением} (scalar multiplication) и записывают следующим образом:

\begin{equation}\label{scalarmultiplication}
(\cdot)\cdot g: \Z_n \to \G\;;\; x \mapsto x\cdot g
\end{equation}
В этой записи символ $x\cdot g$ определяется как «прибавить образующий $g$ к самому себе $x$ раз», а символ $0\cdot g$ определяется как нейтральный элемент в $\G$.
\end{notation}

Криптографические приложения часто используют конечные циклические группы очень большого порядка $n$, что означает, что вычисление экспоненциального отображения путём повторного умножения образующего на самого себя неосуществимо для очень больших классов остатков.\footnote{Однако методы быстрого возведения в степень известны уже давно. Подробное введение можно найти, например, в \chaptname{} 1, \secname{} 7 книги \cite{mignotte-1992}. }
 \term{Алгоритм} (algorithm) \ref{alg_square-and-mul}, называемый \term{возведением в степень квадратами и умножениями} (square and multiply),\sme{unify title of Alg with text} решает эту проблему, вычисляя экспоненциальное отображение примерно за $k$ шагов, где $k$ — битовая длина показателя степени (\ref{def:binary_representation_integer}):\smecomb{SB: I think moving the explanation of bit length here would work better}{move \ref{def:binary_representation_integer} here}

\begin{algorithm}\caption{Возведение в степень в циклической группе}
\label{alg_square-and-mul}
\begin{algorithmic}[0]
\Require $g$ образующий группы порядка $n$
\Require $x \in \Z_n$ 
\Procedure{Exponentiation}{$g,x$}
\State Пусть $(b_0,\ldots,b_k)$ — двоичное представление $x$ \Comment{см. \ref{def:binary_representation_integer}}
\State $h \gets g$
\State $y \gets e_{\G}$
\For{$0\leq j < k$}
	\If{$b_j = 1$}
		\State $y\gets y\cdot h$ \Comment{умножить}
	\EndIf
	\State $h \gets h\cdot h$ \Comment{возвести в квадрат}
\EndFor
\State \textbf{return} $y$
\EndProcedure
\Ensure $ y = g^x$
\end{algorithmic}
\end{algorithm}

Поскольку экспоненциальное отображение уважает групповой закон, не имеет значения, выполняем ли мы наши вычисления в $\Z_n$ до того, как запишем результат в показатель степени $g$, или после: результат в обоих случаях будет одинаковым. Последний метод обычно называют выполнением вычислений «в показателе степени». В криптографии вообще, и в разработке SNARKов в частности, мы часто выполняем вычисления «в показателе степени» образующего.
\begin{example}\label{ex:in-the-exponent} Рассмотрим мультипликативную группу $(\Z_{5}^*,\cdot)$ из упражнения \ref{fstar}. Мы знаем из \exercisename{} \ref{ex:modulus-prime-group}, что $\Z_{5}^*$ — циклическая группа порядка $4$, и что элемент $3\in\Z_5^*$ является образующим. Это означает, что мы также знаем, что следующее отображение уважает групповой закон сложения в $\Z_4$ и групповой закон умножения в $\Z_5^*$:
$$
3^{(\cdot)}: \Z_4 \to \Z_5^* \;;\; x \mapsto 3^x
$$

Чтобы привести пример вычисления «в показателе степени» $3$, выполним вычисление $1+3+2$ в показателе степени образующего $3$:
\begin{align}
3^{1+3+2} &=3^{2}\label{3_Z4_exponent}\\
          & = 4\label{3_Z4_exp_map}
\end{align}
В \eqref{3_Z4_exponent} выше мы сначала выполнили вычисление $1+3+2$ в группе классов остатков $(\Z_4,+)$, а затем применили экспоненциальное отображение $3^{(\cdot)}$ к результату в \eqref{3_Z4_exp_map}. 

Однако поскольку экспоненциальное отображение \eqref{exponentialmap} уважает групповой закон, мы также могли бы отобразить каждое слагаемое в $(\Z_5^*,\cdot)$ сначала, а затем применить групповой закон $(\Z_5^*,\cdot)$. Результат гарантированно будет тем же:

\begin{align*}
3^1 \cdot 3^3 \cdot 3^{2}
          & = 3\cdot 2 \cdot 4\\
          & = 1\cdot 4\\
          & = 4
\end{align*}
\end{example}
Поскольку экспоненциальное отображение \eqref{exponentialmap} является взаимно однозначным соответствием, уважающим групповой закон, можно показать, что это отображение имеет обратное относительно основания $g$, называемое \term{дискретным логарифмическим отображением по основанию g} (base g discrete logarithm map):
\begin{equation}
\label{logarithm_map}
log_g(\cdot): \G \to \Z_n\; x \mapsto log_g(x)
\end{equation}
Дискретные логарифмы чрезвычайно важны в криптографии, поскольку существуют конечные циклические группы, в которых экспоненциальное отображение считается односторонней функцией, что неформально означает, что вычисление экспоненциального отображения быстро, а вычисление логарифмического отображения медленно (мы рассмотрим более точное определение в \ref{crypto_groups}).
\begin{example}Рассмотрим экспоненциальное отображение $3^{(\cdot)}$ из примера \ref{ex:in-the-exponent}. Его обратным является дискретный логарифм по основанию $3$, задаваемый отображением ниже:
$$
log_3(\cdot): \Z_5^* \to \Z_4\; x \mapsto log_3(x)
$$ 
В отличие от экспоненциального отображения $3^{(\cdot)}$, у нас нет способа фактически вычислить это отображение, кроме как перебором всех элементов группы, пока мы не найдём правильный. Например, чтобы вычислить $log_3(4)$, мы должны найти некоторое $x\in \Z_4$ такое, что $3^x=4$, и всё, что мы можем сделать — это многократно подставлять элементы $x$ в показатель степени $3$, пока результат не станет $4$. Чтобы сделать это, запишем все \uterm{образы} $3^{(\cdot)}$: 
$$
\begin{array}{cccc}
3^0 = 1, & 3^1 = 3, & 3^2 = 4, & 3^3 = 2
\end{array}
$$
Поскольку дискретный логарифм $log_3(\cdot)$ определяется как обратная функция к этой функции, мы можем использовать эти образы для вычисления дискретного логарифма:
$$
\begin{array}{ccccc}
log_3(1) = 0, & log_3(2) = 3, & log_3(3) = 1, & log_3(4) = 2
\end{array}
$$
Отметим, что это вычисление было возможно только потому, что мы смогли записать все образы экспоненциального отображения. Однако в реальных приложениях рассматриваемые группы слишком велики, чтобы записывать образы экспоненциального отображения. 
\end{example}
\begin{exercise}[Эффективное скалярное умножение]
\label{alg_double-and-add} Пусть $(\G,+)$ — конечная циклическая группа порядка $n$. Рассмотрите алгоритм \ref{alg_square-and-mul} и определите его аналог для групп в аддитивной записи.
\end{exercise}


\subsection{Факторгруппы}
Как нам известно из основной теоремы арифметики (\ref{def:fundamental_theorem_arithmetic}), каждое натуральное число $n$ является произведением множителей, наиболее простыми из которых являются простые числа. Это параллельно подгруппам конечных циклических групп интересным образом.

\begin{definition}[\deftitle{Основная теорема конечных циклических групп}]\label{def:fundamental_theorem_groups}
Если $\G$ — конечная циклическая группа порядка $n$, то каждая подгруппа $\G'$ группы $\G$ конечна и циклична, а порядок $\G'$ является делителем $n$. Более того, для каждого делителя $k$ числа $n$ группа $\G$ имеет ровно одну подгруппу порядка $k$.
\end{definition}

\begin{notation}Если $\G$ — конечная циклическая группа порядка $n$ и $k$ — делитель $n$, то мы пишем $\G[k]$ для единственной конечной циклической группы, являющейся подгруппой порядка $k$ группы $\G$, и называем её \term{факторгруппой} (factor group) группы $\G$.
\end{notation}

Одна особенно интересная ситуация возникает, если порядок данной конечной циклической группы является простым числом. Как нам известно из основной теоремы арифметики (\ref{def:fundamental_theorem_arithmetic}), простые числа имеют только два делителя: число $1$ и само простое число. Тогда из основной теоремы конечных циклических групп (определение \ref{def:fundamental_theorem_groups}) следует, что такие группы не имеют подгрупп, кроме тривиальной группы (\examplename{} \ref{example:trivial_group}) и самой группы.

Криптографические протоколы часто предполагают существование конечных циклических групп простого порядка. Однако некоторые реальные реализации этих протоколов определены не на группах простого порядка, а на группах, порядок которых состоит из (обычно большого) простого числа, имеющего малые кофакторы (см. \notationname{} \ref{def:cofactor}). В этом случае должен быть применён метод, называемый \term{очисткой кофактора} (cofactor clearing), чтобы гарантировать, что вычисления выполняются не в самой группе, а в её (большой) подгруппе простого порядка.

Чтобы понять очистку кофактора подробно, пусть $\G$ — конечная циклическая группа порядка $n$, и пусть $k$ — делитель $n$ с соответствующей факторгруппой $\G[k]$. Мы можем спроецировать любой элемент $g\in\G[k]$ на нейтральный элемент $e$ группы $\G$, умножив $g$ на самого себя $k$ раз:

\begin{equation}
g^k = e
\end{equation}

Следовательно, если $c:=\Zdiv{n}{k}$ — кофактор $k$ в $n$, то любой элемент из полной группы $g\in \G$ может быть спроецирован в факторгруппу $\G[k]$ умножением $g$ на самого себя $c$ раз. Это определяет следующее отображение, которое в криптографической литературе часто называется \term{очисткой кофактора} (cofactor clearing):
\begin{equation}
\label{def:cofactor_clearing}
(\cdot)^c: \G \to \G[k]\; : \; g \mapsto g^c
\end{equation}

\begin{example}\label{example:factor_groupds_of_F*5} Рассмотрим конечную циклическую группу $(\Z_5^*,\cdot)$ из примера \ref{example:cyclic_group_F5*}. Поскольку порядок $\Z^*_5$ равен $4$, а $4$ имеет делители $1$, $2$ и $4$, из основной теоремы конечных циклических групп (\defname{} \ref{def:fundamental_theorem_groups}) следует, что $\Z^*_5$ имеет $3$ различные подгруппы. На самом деле, единственная подгруппа $\Z^*_5[1]$ порядка $1$ задаётся тривиальной группой $\{1\}$, содержащей только мультипликативный нейтральный элемент $1$. Единственная подгруппа $\Z^*_5[4]$ порядка $4$ есть сама $\Z^*_5$, так как по определению каждая группа тривиально является подгруппой самой себя. Единственная подгруппа $\Z^*_5[2]$ порядка $2$ более интересна и задаётся множеством $\Z^*_5[2]=\{1,4\}$.

Поскольку $\Z^*_5$ не является группой простого порядка, и поскольку единственный простой делитель $4$ — это $2$, «большая» подгруппа простого порядка $\Z^*_5$ есть $\Z^*_5[2]$. Кроме того, поскольку кофактор $2$ в $4$ также равен $2$, мы получаем отображение очистки кофактора $(\cdot)^2:\Z^*_5 \to \Z^*_5[2]$. Как и ожидалось, когда мы применяем это отображение ко всем элементам $\Z^*_5$, мы видим, что оно отображает только на элементы $\Z^*_5[2]$:

\begin{equation}
1^2 = 1\quad{}  2^2 = 4\quad{}  3^2 = 4\quad{}  4^2 = 1
\end{equation}

Мы можем использовать это отображение, чтобы «очистить кофактор» любого элемента из $\Z^*_5$, что означает, что элемент спроецирован на «большую» подгруппу простого порядка $\Z^*_5[2]$.
\end{example}
\begin{exercise}Рассмотрите предыдущий \examplename{} \ref{example:factor_groupds_of_F*5} и покажите, что $\Z^*_5[2]$ является коммутативной группой.
\end{exercise}
\begin{exercise}
Рассмотрим конечную циклическую группу $(\Z_6,+)$ сложения по модулю $6$ из \examplename{} \ref{example:cyclic_group_F6}. Опишите все подгруппы $(\Z_6,+)$. Определите большую подгруппу простого порядка $\Z_6$, определите её отображение очистки кофактора и примените это отображение ко всем элементам $\Z_6$.
\end{exercise}
\begin{exercise}Пусть $(\Z_p^*,\cdot)$ — циклическая группа из \exercisename{} \ref{ex:modulus-prime-group}. Покажите, что для $p\geq 5$ не каждый элемент $x\in \F_p^*$ является образующим $\F_p^*$.
\end{exercise}
\subsection{Паринги}
Особого значения для разработки SNARKов имеют так называемые \term{паринговые отображения} (pairing maps) на коммутативных группах, определённые ниже.

\begin{definition}[\deftitle{Паринговое отображение}]\label{def:pairing-map} 
Пусть $\G_1$, $\G_2$ и $\G_3$ — три коммутативные группы. Тогда \textbf{паринговое отображение} — это функция
\begin{equation}\label{pairing-map}
e(\cdot,\cdot): \G_1 \times \G_2 \to \G_3
\end{equation}

Эта функция принимает пары $(g_1,g_2)$ элементов из $\G_1$ и $\G_2$ и отображает их в элементы $\G_3$ так, что выполняется свойство \term{билинейности} (bilinearity), то есть для всех $g_1,g_1'\in \G_1$ и $g_2, g_2'\in \G_2$ выполняются следующие два тождества:
\begin{equation}\label{eq:bilinearity}
\begin{array}{lcr}
e(g_1 \cdot g_1',g_2)= e(g_1,g_2)\cdot e(g_1',g_2) &\text{и}&
e(g_1,g_2 \cdot g_2')= e(g_1,g_2)\cdot e(g_1,g_2')\\
\end{array}
\end{equation}
\end{definition}

Неформально говоря, билинейность означает, что не имеет значения, сначала ли выполняем групповой закон на одной стороне, а затем применяем билинейное отображение, или сначала применяем билинейное отображение, а затем групповой закон в $\G_3$. 

Паринговое отображение называется \term{невырожденным} (non-degenerate), если всякий раз, когда результат паринга является нейтральным элементом в $\G_3$, одно из входных значений является нейтральным элементом $\G_1$ или $\G_2$. Более точно, $e(g_1,g_2)=e_{\G_3}$ влечёт $g_1=e_{\G_1}$ или $g_2=e_{\G_2}$.

\begin{example}
\label{example:integer_addition_pairing}
Один из самых простых примеров невырожденного паринга включает группы $\G_1$, $\G_2$ и $\G_3$, все являющиеся группами целых чисел со сложением $(\Z,+)$. В этом случае следующее отображение определяет невырожденный паринг:
\begin{equation}
e(\cdot,\cdot): \Z \times \Z \to \Z \; (a,b)\mapsto a\cdot b
\end{equation}

Отметим, что билинейность следует из дистрибутивного закона целых чисел, то есть для $a,b,c\in \Z$ выполняется равенство $e(a+b,c)=(a+b)\cdot c = a\cdot c + b\cdot c = e(a,c)+ e(b,c)$ (и то же рассуждение верно для второго аргумента \smecomb{$b$}{$b$}).

Чтобы увидеть, что $e(\cdot,\cdot)$ невырождено, предположим, что $e(a,b)=0$. Тогда $a\cdot b =0$ влечёт, что $a$ или $b$ должны быть равны нулю.
\end{example}
\begin{exercise}[Арифметические законы для паринговых отображений]
\label{ex:pairing-arithmetics} Пусть $\G_1$, $\G_2$ и $\G_3$ — конечные циклические группы того же порядка $n$, и пусть $e(\cdot,\cdot): \G_1 \times \G_2 \to \G_3$ — паринговое отображение. Покажите, что для данных $g_1\in \G_1$, $g_2\in \G_2$ и всех $a,b\in \Z_n$ выполняется следующее тождество:
\begin{equation}
e(g_1^a, g_2^b) = e(g_1,g_2)^{a\cdot b}
\end{equation}
\end{exercise}
\begin{exercise} Рассмотрим группы классов остатков $(\Z_n,+)$ из \examplename{} \ref{example:Zn} для некоторого модуля $n$. Покажите, что следующее отображение является паринговым отображением. 
\begin{equation}
e(\cdot,\cdot): \Z_n \times \Z_n \to \Z_n \; (a,b)\mapsto a\cdot b
\end{equation}

Почему паринг в общем случае не невырожден, и какое условие должно быть наложено на $n$, чтобы паринг был невырожденным?
\end{exercise}

\subsection{Криптографические группы}
\label{crypto_groups} В этом разделе мы рассмотрим классы групп, которые, как считается, удовлетворяют определённым \term{предположениям вычислительной трудности} (computational hardness assumptions), то есть вычисление которых за полиномиальное время невозможно.\footnote{Более подробное введение в предположения вычислительной трудности и их приложения в криптографии можно найти в \chaptname{} 3, \secname{} 8 книги \cite{katz-2007}.}

\begin{example}\label{ex:computational-hardness}

Чтобы привести пример хорошо известного предположения вычислительной трудности, рассмотрим задачу факторизации, то есть вычисления простых множителей составного целого числа (см. \examplename{} \ref{ex-prime-factorization}). Если простые множители очень велики, это невозможно сделать, и ожидается, что так оно и останется. Мы предполагаем, что задача \term{вычислительно трудна} (computationally hard) или \term{невыполнима} (infeasible).

\end{example}

Отметим, что в \examplename{} \ref{ex:computational-hardness} {мы говорим, что} задача невыполнима для решения \hilight{если простые множители достаточно велики}. Естественно, это уточняется в криптографической стандартной модели, где у нас есть параметр безопасности, и мы говорим, что «существует параметр безопасности такой, что невозможно вычислить решение задачи». В следующих примерах параметр безопасности примерно коррелирует с порядком рассматриваемой группы. В этой книге мы не включаем параметр безопасности в наши определения, поскольку наша цель — только дать интуитивное понимание криптографических предположений, а не научить способности проводить строгий анализ.

Кроме того, поймите, что это \hilight{предположения}. Академики давно ищут эффективные алгоритмы разложения на простые множители, и они становятся всё лучше и лучше, а компьютеры — всё быстрее и быстрее — но всегда находился более высокий параметр безопасности, для которого задача всё ещё оставалась невыполнимой.% Это также объясняет, почему новые ключи RSA имеют длину 4096 бит, в отличие от 1024 бит 20 лет назад.

В дальнейшем мы опишем несколько задач, возникающих в контексте групп в криптографии, которые предполагаются невыполнимыми. Мы будем ссылаться на них на протяжении всей книги.


\subsubsection{Проблема \capitalisewords{дискретного логарифма}}
\label{def:DL-secure}
Так называемая \term{\capitalisewords{проблема дискретного логарифма} (DLP)} (discrete logarithm problem), также называемая \term{\capitalisewords{предположением дискретного логарифма}} (discrete logarithm assumption), является одним из самых фундаментальных предположений в криптографии. 

\begin{definition}[]
Пусть $\G$ — конечная циклическая группа порядка $r$ и $g$ — образующий $\G$. Мы знаем из \eqref{exponentialmap}, что существует экспоненциальное отображение $g^{(\cdot)}: \Z_r \to \G\; ;\; x\mapsto g^x$, которое отображает классы вычетов арифметики по модулю $r$ на группу во взаимно однозначном соответствии.
\textbf{\capitalisewords{Проблема дискретного логарифма}} — это задача нахождения обратного к этому отображению, то есть нахождения решения $x\in\Z_r$ следующего уравнения для некоторых данных $h, g \in \G$:

\begin{equation}
h = g^x
\end{equation}
\end{definition}

Существуют группы, в которых предполагается, что DLP невозможно решить, и есть группы, в которых это не так. Мы называем первые \term{DL-безопасными} (DL-secure) группами.

Перефразируя предыдущее определение, считается, что в DL-безопасных группах существует число $n$ такое, что невозможно вычислить некоторое число $x$, которое решает уравнение $h=g^x$ для данных $h$ и $g$, при условии, что порядок группы $n$ достаточно велик. Число $n$ здесь соответствует параметру безопасности, обсуждённому выше.

\begin{example}[Криптография с открытым ключом]\label{ex:publ-key}

Один из самых базовых примеров применения DL-безопасных групп — криптография с открытым ключом, где стороны публично договариваются о некоторой паре $(\G,g)$ такой, что $\G$ — конечная циклическая группа соответствующего порядка $n$, предположительно являющаяся DL-безопасной группой, а $g$ — образующий $\G$.

В этой постановке секретный ключ — это некоторое число $sk \in \Z_r$, а соответствующий открытый ключ $pk$ — элемент группы $pk=g^{sk}$. Поскольку дискретные логарифмы предполагаются трудными, для злоумышленника невозможно вычислить секретный ключ из открытого ключа, так как это потребовало бы нахождения решений $x$ следующего уравнения (что предполагается невыполнимым):

\begin{equation}
pk = g^{x}
\end{equation}

\end{example}

Как показывает \examplename{} \ref{ex:publ-key}, идентификация DL-безопасных групп — важная практическая задача. К сожалению, легко видеть, что предполагать трудность \capitalisewords{проблемы дискретного логарифма} во всех конечных циклических группах не имеет смысла: контрпримеры распространены и легко построить.\sme{mention a few examples}

\begin{comment}
\begin{example}[Модульная арифметика для простых чисел Ферма]

Широко считается, что \capitalisewords{проблема дискретного логарифма} трудна в мультипликативных группах $\Z_p^*$ арифметики по модулю простого числа. Однако не все такие группы являются DL-безопасными. Чтобы увидеть это, рассмотрим любое так называемое простое число Ферма, то есть простое число $p\in\Prim$ такое, что $p=2^n+1$ для некоторого числа $n$.

Мы знаем из упражнения \ref{ex:Zn*}\sme{check reference}, что в этом случае $\Z_p^* = \{1,2,\ldots, p-1\}$ является группой относительно умножения целых чисел в арифметике по модулю $p$, и поскольку $p=2^n+1$, порядок $\Z_p^*$ равен $2^n$, что означает, что соответствующий параметр безопасности задаётся $log_2(2^n)=n$.

Мы показываем, что в этом случае $\Z_p^*$ не является DL-безопасной группой, построив алгоритм, способный вычислить некоторое $x\in\Z_{2^n}$ для любого данных образующего $g$ и произвольного элемента $h$ из $\F_p^*$ так, что выполняется уравнение \ref{eq:hgx}, и вычислительная сложность построенного алгоритма составляет $\mathcal{O}(n^2)$, что квадратично относительно параметра безопасности $n=log_2(2^n)$.

\begin{equation}
h = g^x
\end{equation}

(eq:hgx)

Чтобы определить такой алгоритм, предположим, что образующий $g$ является публичной константой и дан элемент группы $h$. Наша задача — эффективно вычислить $x$.

Первое, на что следует обратить внимание, заключается в том, что поскольку $x$ — число в арифметике по модулю $2^n$, мы можем записать двоичное представление $x$, как в \ref{eq:binary-x}, с двоичными коэффициентами $c_j\in\{0,1\}$. В частности, $x$ — это $n$-битное число, если интерпретировать его как целое.\sme{explain last sentence more}

\begin{equation}
x = c_0\cdot 2^0 + c_1\cdot 2^1 + \cdots + c_n \cdot 2^n
\end{equation}

(eq:binary-x)


Мы затем используем это представление для построения алгоритма, который вычисляет биты $c_j$ один за другим, начиная с $c_0$. Чтобы понять, как этого можно достичь, заметим, что мы можем определить $c_0$, возведя вход $h$ в степень $2^{n-1}$ в $\F_p^*$. Мы используем законы степеней и вычисляем следующим образом:
\begin{align*}
h^{2^{n-1}} & = \left(g^x\right)^{2^{n-1}}\\
            & = \left(g^{c_0\cdot 2^0 + c_1\cdot 2^1 + \ldots + c_n\cdot 2^n}\right)^{2^{n-1}}\\
            & = g^{c_0\cdot 2^{n-1}}\cdot g^{c_1\cdot 2^1\cdot 2^{n-1}} \cdot
            g^{c_2\cdot 2^2\cdot 2^{n-1}} \cdots g^{c_n\cdot 2^n\cdot 2^{n-1}}\\
            & = g^{c_0 2^{n-1}}\cdot g^{c_1\cdot 2^0\cdot 2^{n}} \cdot
            g^{c_2\cdot 2^1\cdot 2^{n}} \cdots g^{c_n\cdot 2^{n-1}\cdot 2^{n}}
\end{align*}
Теперь, поскольку $g$ — образующий и $\F_p^*$ циклична порядка $2^n$, мы знаем $g^{2^n}=1$ и поэтому $g^{k\cdot 2^n}= 1^k=1$. Отсюда следует, что все сомножители, кроме первого, в последнем выражении равны $1$, и мы можем упростить выражение до следующего:
\begin{equation}
h^{2^{n-1}} = g^{c_0 2^{n-1}}
\end{equation}
Теперь в случае $c_0=0$ мы получаем $h^{2^{n-1}} = g^0=1$. В случае $c_0=1$ мы получаем
$h^{2^{n-1}} = g^{2^{n-1}}\neq 1$ (Чтобы увидеть, что $g^{2^{n-1}}\neq 1$, вспомним, что $g$ — образующий $\F_p^*$, и следовательно, $\F_p^*$ — циклическая группа порядка $2^n$, что влечёт $g^y\neq 1$ для всех $y<2^n$).

Возведение $h$ в степень $2^{n-1}$ определяет $c_0$, и мы можем применить то же рассуждение к коэффициенту $c_1$, возведя $h\cdot g^{-c_0\cdot 2^0}$ в степень $2^{n-2}$. Этот подход затем может быть повторён, пока все коэффициенты $c_j$ числа $x$ не будут найдены.

Предполагая, что возведение в степень в $\F_p^*$ может быть выполнено с логарифмической вычислительной сложностью $log(p)$, следует, что наш алгоритм имеет вычислительную сложность
$\mathcal{O}(log^2(p))=\mathcal{O}(n^2)$, поскольку мы должны выполнить $n$ возведений в степень, чтобы определить $n$ двоичных коэффициентов $x$.

Отсюда следует, что всякий раз, когда $p$ — простое число Ферма, предположение дискретного логарифма не выполняется в $F_p^*$.

\end{example}
\end{comment}

\subsubsection{Предположение \capitalisewords{диффи — хеллмана о различении}}
\label{def:DDH-secure}

\begin{definition}
Пусть $\G$ — конечная циклическая группа порядка $n$ и $g$ — образующий $\G$. Проблема диффи — хеллмана о различении (DDH) заключается в различении $(g^a,g^b, g^{ab})$ от тройки $(g^a,g^b,g^c)$ для равномерно случайных значений $a,b,c\in \Z_r$. 
\end{definition}

Если мы предполагаем, что проблема DDH невозможна для решения в $\G$, мы называем $\G$ \term{DDH-безопасной} (DDH-secure) группой.

DDH-безопасность — более сильное предположение, чем DL-безопасность (\ref{def:DL-secure}), в том смысле, что если проблема DDH невыполнима, то и DLP невыполним, но не обязательно наоборот.

Чтобы понять, почему это так, предположим, что предположение дискретного логарифма не выполняется. В этом случае, данные образующий $g$ и элемент группы $h$, легко вычислить некоторый элемент $x\in\Z_p$ с $h=g^x$. Тогда предположение диффи — хеллмана о различении не может выполняться, так как данная тройка $(g^a , g^b , z )$, можно эффективно решить, верно ли $z = g^{ab}$, сначала вычислив дискретный логарифм $b$ числа $g^b$, затем вычислив $g^{ab}= (g^a)^b$ и решив, верно ли $z=g^{ab}$.

С другой стороны, следующий пример показывает, что существуют группы, где предположение дискретного логарифма выполняется, а \capitalisewords{предположение диффи — хеллмана о различении} — нет.

\begin{example}[Эффективно вычислимые билинейные паринги]

Пусть $\G$ — DL-безопасная, конечная, циклическая группа порядка $r$ с образующим $g$, а $\G_T$ — другая группа такая, что существует эффективно вычислимое паринговое отображение $e(\cdot,\cdot): \G \times \G \to \G_T$, которое билинейно и невырождено (\ref{pairing-map}).

В такой постановке легко показать, что решение DDH не может быть невыполнимым, так как, данная тройка $(g^a, g^b, z)$, возможно эффективно проверить, верно ли $z = g^{ab}$, используя следующий паринг:

\begin{equation}
e(g^a,g^b) \checkeq e(g,z)
\end{equation}

Поскольку свойства билинейности $e(\cdot,\cdot)$ влекут $e(g^a,g^b)= e(g,g)^{ab}= e(g,g^{ab})$, а $e(g,y)=e(g,y')$ влечёт $y=y'$ из-за невырожденности, равенство означает $z=g^{ab}$.

\end{example}

Следовательно, предположение DDH действительно сильнее предположения дискретного логарифма, и группы с эффективными парингами не могут быть DDH-безопасными группами. %Следующий пример показывает ещё один важный класс групп, в которых DDH-безопасность не выполняется: мультипликативные группы классов вычетов по простому модулю.

\begin{comment}
\begin{example}

Пусть $p$ — простое число и $\F_p^*=\{1,2,\ldots,p-1\}$ — мультипликативная группа арифметики по модулю $p$, как в упражнении \ref{ex:Zn*}. Чтобы увидеть, что $\F_p^*$ не может быть DDH-безопасной группой, вспомним из XXX, что \uterm{символ Лежандра} $\legsym{x}{p}$ любого $x\in \F_p^*$ эффективно вычисляется по \uterm{формуле Эйлера}.\sme{These are only explained later in the text, `\ref{eq: Legendre-symbol}`} Но символ Лежандра числа $g^{a}$ раскрывает, чётно ли $a$. Даны $g^{a}$, $g^{b}$ и $g^{ab}$, можно таким образом эффективно вычислить и сравнить наименее значащие биты $a$, $b$ и $ab$ соответственно, что даёт вероятностный метод отличить $g^{ab}$ от случайного элемента группы $g^c$.

\end{example}
\end{comment}

\subsubsection{Предположение \capitalisewords{вычислительной сложности диффи — хеллмана}}
%
\begin{definition}
Пусть $\G$ — конечная циклическая группа порядка $n$ и $g$ — образующий $\G$. \term{Предположение вычислительной сложности диффи — хеллмана} (computational Diffie--Hellman assumption) предполагает, что данные случайно и независимо выбранные элементы $a,b\in\Z_r$, невозможно вычислить $g^{ab}$, если известны только $g$, $g^a$ и $g^b$ (но не $a$ и $b$). Если это так для $\G$, мы называем $\G$ \term{CDH-безопасной} (CDH-secure) группой.
\end{definition}

Вообще говоря, мы не знаем, является ли CDH-безопасность более сильным предположением, чем DL-безопасность, или оба предположения эквивалентны. Мы знаем, что DL-безопасность необходима для CDH-безопасности, но обратное направление в настоящее время недостаточно изучено. В частности, неизвестны DL-безопасные группы, которые не были бы также CDH-безопасными. %\citep{Fifield12theequivalence}.% https://web.stanford.edu/class/cs259c/finalpapers/dlp-cdh.pdf

Чтобы понять, почему предположение дискретного логарифма необходимо, предположим, что оно не выполняется. Тогда, данные образующий $g$ и элемент группы $h$, легко вычислить некоторый элемент $x\in\Z_p$ с $h=g^x$. В этом случае предположение вычислительной сложности диффи — хеллмана не может выполняться, так как, даны $g$, $g^a$ и $g^b$, возможно эффективно вычислить $b$, то есть $g^{ab}=(g^a)^b$ может быть вычислено из этих данных.

Предположение вычислительной сложности диффи — хеллмана — более слабое предположение, чем \capitalisewords{предположение диффи — хеллмана о различении}. Это означает, что существуют группы, где CDH выполняется, а DDH — нет, тогда как не может быть групп, в которых DDH выполняется, а CDH — нет. Чтобы увидеть это, предположим, что возможно эффективно вычислить $g^{ab}$ из $g$, $g^a$ и $g^b$. Тогда, дана $(g^a,g^b,z)$, легко решить, выполняется ли $z=g^{ab}$ или нет.

Существует несколько вариаций и частных случаев CDH. Например, \term{предположение \capitalisewords{квадратной вычислительной сложности диффи — хеллмана}} (square computational Diffie--Hellman assumption) предполагает, что, даны $g$ и $g^x$, вычислительно трудно вычислить $g^{x^2}$. \term{Предположение \capitalisewords{обратной вычислительной сложности диффи — хеллмана}} (inverse computational Diffie--Hellman assumption) предполагает, что, даны $g$ и $g^x$, вычислительно трудно вычислить $g^{x^{-1}}$.

\subsection{Хеширование в группы}\label{sec:hashing-to-groups}
% https://crypto.stackexchange.com/questions/78017/simple-hash-into-a-prime-field
\subsubsection{Хеш-функции}\label{sec:hash-functions} Вообще говоря, хеш-функция — это любая функция, которая может использоваться для отображения данных произвольного размера в значения фиксированного размера. Поскольку двоичные строки произвольной длины — это способ представления данных вообще, мы можем понимать \textbf{хеш-функцию} как следующее отображение, где $\{0,1\}^*$ представляет множество всех двоичных строк произвольной, но конечной длины, а $\{0,1\}^k$ представляет множество всех двоичных строк, имеющих длину ровно $k$ бит:
\begin{equation}
\label{def:hash_function}
H: \{0,1\}^* \to \{0,1\}^k
\end{equation}
\term{Образами} (images) $H$, то есть значениями, возвращаемыми хеш-функцией $H$, называются \term{хеш-значениями} (hash values), \term{дайджестами} (digests) или просто \term{хешами} (hashes).

\begin{notation}
\label{string_and_hash_notations}
В дальнейшем мы называем элемент $b\in\{0,1\}$ \term{битом} (bit). Если $s\in\{0,1\}^*$\sme{check footnote} — двоичная строка, мы пишем $|s|=k$ для её \term{длины} (length), то есть для числа битов в $s$. Мы пишем $<>$ для пустой двоичной строки и $s=<b_1,b_2,\ldots,b_k>$ для двоичной строки длины $k$.\footnote{Разница между обозначениями $b\in\{0,1\}$ и $s\in\{0,1\}^*$ следующая: $b\in\{0,1\}$ означает, что $b$ равно либо $0$, либо $1$, тогда как $s$ — это строка, состоящая из произвольного числа $0$ и $1$ (и $s$ также может быть пустой строкой).}

Если даны две двоичные строки $s=<b_1,b_2,\ldots,b_k>$ и $s'=<b'_1,b'_2,\ldots,b'_l>$, то мы пишем $s||s'$ для \term{конкатенации} (concatenation), то есть строки 
$s||s'=<b_1,b_2,\ldots,b_k,b'_1,b'_2,\ldots,b'_l>$.

Если $H$ — хеш-функция, отображающая двоичные строки произвольной длины в двоичные строки длины $k$, а $s\in\{0,1\}^*$ — двоичная строка, мы пишем $H(s)_j$ для бита в позиции $j$ в образе $H(s)$.
\end{notation}

\begin{example}[Хеш с $k$-усечением]\label{ex:k-truncation-hash} Одна из самых простых хеш-функций $H_k:\{0,1\}^*\to \{0,1\}^k$ задаётся простым усечением каждой двоичной строки $s$ размера $|s|> k$ до строки размера $k$ и дополнением любой строки $s'$ размера $|s'|<k$ нулями. Чтобы сделать эту хеш-функцию детерминированной, мы определяем, что и усечение, и дополнение должны происходить с наиболее значимых битов, или «слева».

Например, если параметр $k$ равен $k=3$, $s_1=<0,0,0,0,1,0,1,0,1,1,1,0>$ и $s_2=<1>$, то $H_3(s_1)=<1,1,0>$ и $H_3(s_2)=<0,0,1>$.
\end{example}

Желательным свойством хеш-функции является \term{равномерность} (uniformity), то есть она должна отображать входные значения как можно более равномерно по диапазону выходных значений. В математических терминах, каждая строка длины $k$ из $\{0,1\}^k$ должна генерироваться примерно с одинаковой вероятностью.

Особого интереса являются так называемые \term{криптографические} (cryptographic) хеш-функции, которые являются хеш-функциями, являющимися также \term{односторонними функциями} (one-way functions), что по существу означает, что, дана строка $y$ из $\{0,1\}^k$, невозможно найти строку $x\in\{0,1\}^*$ такую, что $H(x)=y$. Это свойство обычно называется \term{стойкостью к нахождению прообраза} (preimage resistance).

Кроме того, если дана строка $x_1\in\{0,1\}^*$, то должно быть невозможно найти другую строку $x_2\in\{0,1\}^*$ с $x_1\neq x_2$ и $H(x_1)=H(x_2)$. Это часто называется свойством \term{стойкости к нахождению второго прообраза} (second preimage resistance).

В дополнение, должно быть невозможно найти две различные строки $x_1,x_2 \in\{0,1\}^*$ такие, что $H(x_1)=H(x_2)$, что называется \term{стойкостью к коллизиям} (collision resistance). Важно отметить, однако, что коллизии всегда существуют, поскольку функция $H: \{0,1\}^* \to \{0,1\}^k$ неизбежно отображает бесконечно много значений на один и тот же хеш. На самом деле, для любой хеш-функции с дайджестами длины $k$ нахождение прообраза для данного дайджеста всегда может быть выполнено полным перебором за $2^k$ шагов оценки. Просто должно быть практически невозможно вычислить эти значения, и статистически очень маловероятно сгенерировать два из них случайно.

Третьим свойством криптографической хеш-функции является то, что малые изменения во входной строке, например изменение одного бита, должны генерировать хеш-значения, которые выглядят совершенно по-разному. Это называется \term{диффузией} (diffusion) или лавинным эффектом (avalanche effect).

Поскольку криптографические хеш-функции отображают крошечные изменения входных значений на большие изменения в выходных, ошибки реализации, изменяющие результат, обычно легко заметить, сравнивая их с ожидаемыми выходными значениями. Поэтому определения криптографических хеш-функций обычно сопровождаются некоторыми тестовыми векторами распространённых входных данных и ожидаемых дайджестов. Поскольку пустая строка $<>$ — единственная строка длины $0$, распространённым тестовым вектором является ожидаемый дайджест пустой строки.
\begin{example}[Хеш с $k$-усечением] Рассмотрим хеш с $k$-усечением из примера \ref{ex:k-truncation-hash}. Поскольку пустая строка имеет длину $0$, следует, что дайджест пустой строки — это строка длины $k$, содержащая только $0$:
\begin{equation}
H_k(<>)= <0,0,\ldots, 0,0>
\end{equation}
Вполне очевидно из определения $H_k$, что эта простая хеш-функция не является криптографической хеш-функцией. В частности, каждый дайджест является собственным прообразом, поскольку $H_k(y)=y$ для каждой строки размера ровно $k$. Поэтому нахождение прообразов легко, и свойство стойкости к нахождению прообраза не выполняется.

Кроме того, легко построить коллизии, так как все строки $s$ размера $|s|>k$, имеющие одинаковые $k$ битов «справа», отображаются на одно и то же хеш-значение. Это означает, что эта функция также не является стойкой к коллизиям.

Наконец, у этой хеш-функции небольшая диффузия, так как изменение битов, не входящих в $k$ крайних правых битов, вообще не изменяет дайджест.
\end{example}

Вычисление криптографически стойких хеш-функций в стиле «бумага и ручка» возможно, но утомительно. К счастью, Sage может импортировать библиотеку \term{\code{hashlib}}, которая предназначена для обеспечения надёжной и стабильной основы для написания программ Python, требующих криптографических функций. Следующие примеры объясняют, как использовать \code{hashlib} в Sage.

\begin{example}\label{ex:SHA256}Примером хеш-функции, которая обычно считается криптографически стойкой, является так называемая \term{SHA256}, которая в наших обозначениях является функцией, отображающей двоичные строки произвольной длины в двоичные строки длины $256$:
\begin{equation}
SHA256: \{0,1\}^* \to \{0,1\}^{256}
\end{equation}

 Чтобы оценить правильную реализацию хеш-функции $SHA256$, дайджест пустой строки должен быть следующим:
 
\begin{equation}
SHA256(<>)= {\scriptstyle e3b0c44298fc1c149afbf4c8996fb92427ae41e4649b934ca495991b7852b855}
\end{equation}

Для лучшей читаемости человеком распространённой практикой является представление дайджеста строки не в двоичной форме, а в шестнадцатеричном представлении. Мы можем использовать Sage для вычисления SHA256 и свободно переходить между двоичным, шестнадцатеричным и десятичным представлениями. Для этого мы импортируем реализацию SHA256 из \code{hashlib}:
\begin{sagecommandline}
sage: import hashlib
sage: test = 'e3b0c44298fc1c149afbf4c8996fb92427ae41e4649b934ca495991b7852b855' 
sage: empty_string = ""
sage: binary_string = empty_string.encode()
sage: hasher = hashlib.sha256(binary_string) 
sage: result = hasher.hexdigest()
sage: type(result)	# Sage представляет дайджесты как строки
sage: d = ZZ('0x'+ result) # преобразование в целое число
sage: d.str(16) == test	# хеш равен тестовому вектору
sage: d.str(16) # шестнадцатеричное представление
sage: d.str(2) # двоичное представление
sage: d.str(10) # десятичное представление
\end{sagecommandline}
\end{example}

\subsubsection{Хеширование в циклические группы} Как мы видели в предыдущем разделе, общие хеш-функции отображают двоичные строки произвольной длины в двоичные строки некоторой фиксированной длины. Однако в различных криптографических примитивах желательно не просто хешировать в двоичные строки фиксированной длины, а хешировать в алгебраические структуры, такие как группы, сохраняя при этом (некоторые) свойства хеш-функции, такие как стойкость к нахождению прообраза или стойкость к коллизиям.

Хеш-функции такого рода могут быть определены для различных алгебраических структур, но, в некотором смысле, наиболее фундаментальными являются хеш-функции, отображающие в группы, поскольку они легко могут быть расширены для отображения в другие структуры, такие как кольца или поля.

Чтобы дать более точное определение, пусть $\G$ -- группа и $\{0,1\}^*$ -- множество всех конечных двоичных строк, тогда \term{хеш-в-группу} (hash-to-group) функция -- это детерминированное отображение:
\begin{equation}
H : \{0,1\}^* \to \G
\end{equation}

Как показывает следующий пример, хеширование в конечные циклические группы может быть тривиально достигнуто ценой некоторых нежелательных свойств хеш-функции:

\begin{example}[Наивное хеширование в циклическую группу]\label{naive-cyclic-group-hash} Пусть $\G$ -- конечная циклическая группа порядка $n$. Если задача состоит в реализации хеш-в-группу функции, один непосредственный подход может основываться на наблюдении, что двоичные строки размера $k$ могут быть интерпретированы как целые числа $z\in\Z$ в диапазоне $0\leq z < 2^k$ с использованием уравнения \ref{def:binary_representation_integer}.

Более точно, пусть $H:\{{0,1}\}^*\to \{0,1\}^k$ -- хеш-функция для некоторого параметра $k$, $g$ -- образующий $\G$, а $s\in\{0,1\}^*$ -- двоичная строка. Используя уравнение \ref{def:binary_representation_integer} и обозначение \ref{string_and_hash_notations}, следующее выражение является неотрицательным целым числом:

\begin{equation}
z_{H(s)}= H(s)_0\cdot 2^0 + H(s)_1\cdot 2^1 + \ldots + H(s)_k \cdot 2^k
\end{equation}

Хеш-в-группу функция для группы $\G$ может быть определена как композиция экспоненциального отображения $g^{(\cdot)}$ числа $g$ с интерпретацией $H(s)$ как целого числа:

\begin{equation}
H_{g} : \{0,1\}^* \to \G:\; s \mapsto g^{z_{H(s)}}
\end{equation}

Построение хеш-в-группу функции подобным образом легко для циклических групп, и оно может быть достаточно хорошим в определённых приложениях.\sme{a few examples?} Однако оно почти никогда не является адекватным в криптографических приложениях, поскольку дискретно-логарифмическое соотношение может быть построено между некоторыми хеш-значениями $H_g(s)$ и $H_g(t)$, независимо от того, является ли $\G$ DL-безопасной (см. \secname \ref{def:DL-secure}).

Более точно, дискретно-логарифмическое соотношение между элементами группы $H_g(s)$ и $H_g(t)$ -- это любой элемент $x\in \Z_n$ такой, что $H_g(s) = H_g(t)^x$. Чтобы понять, как такое $x$ может быть построено, предположим, что $z_{H(s)}$ имеет мультипликативный обратный в $\Z_n$. В этом случае элемент $x=z_{H(t)}\cdot z_{H(s)}^{-1}$ из $\Z_n$ является дискретно-логарифмическим соотношением между $H_g(s)$ и $H_g(t)$:
\begin{align*}
g^{z_{H(t)}} & = g^{z_{H(t)}} & \Leftrightarrow\\
g^{z_{H(t)}} & = g^{z_{H(t)}\cdot z_{H(s)}\cdot z_{H(s)}^{-1}} & \Leftrightarrow \\
g^{z_{H(t)}} & = g^{z_{H(s)}\cdot x} & \Leftrightarrow \\
H_g(t) & = (H_g(s))^x
\end{align*}
\end{example}
Поэтому приложениям, в которых дискретно-логарифмические соотношения между хеш-значениями нежелательны, нужны другие подходы. Многие из этих подходов начинаются со способа хеширования в множество $\Z_r$ арифметики по модулю $r$.

\subsubsection{\capitalisewords{Хеши Педерсена}}
\label{def:Pedersen_hash}
Так называемая \term{\capitalisewords{хеш-функция Педерсена}} (Pedersen hash function) \citep{Pedersen92} предоставляет способ отображения кортежей фиксированного размера элементов из модульной арифметики на элементы конечных циклических групп таким образом, что дискретно-логарифмические соотношения (см. \examplename{} \ref{naive-cyclic-group-hash}) между различными образами исключаются. Композиции \capitalisewords{хеша Педерсена} с общей хеш-функцией \eqref{def:hash_function} затем предоставляют хеш-в-группу функции, отображающие строки произвольной длины на элементы групп.

Более точно, пусть $j$ -- целое число, $\G$ -- конечная циклическая группа порядка $n$, а $\{g_1, \ldots, g_j\} \subset \G$ -- равномерно и случайно сгенерированное множество образующих $\G$. Тогда \term{хеш-функция Педерсена} (Pedersen's hash function) определяется следующим образом:
\begin{equation}\label{eq:Pedersen_hash}
H_{\{g_1,\ldots,g_j\}} : \left(\Z_r\right)^j \to \G:\; (x_1,\ldots,x_j)\mapsto \Pi_{i=1}^j g_i^{x_i}
\end{equation}

Можно показать, что хеш-функция Педерсена стойка к коллизиям при предположении, что $\G$ является DL-безопасной (см. \secname \ref{def:DL-secure}). Важно отметить, однако, что следующее семейство функций не квалифицируется как \uterm{семейство псевдослучайных функций} (pseudorandom function family).

\begin{equation}
\label{Pedersen_not_pseudorandom}
\{H_{\{g_1,\ldots,g_j\}}\;|\;g_1,\ldots,g_j\in \G\}
\end{equation}

С точки зрения реализации, важно выводить множество образующих $\{g_1,\ldots,g_k\}$ таким образом, чтобы они были как можно более равномерными и случайными. В частности, должно быть исключено любое известное дискретно-логарифмическое соотношение между двумя образующими, то есть любое известное $x\in \Z_n$ с $g_h = (g_i)^x$.

\begin{example} Чтобы вычислить фактический хеш Педерсена, рассмотрим циклическую группу $\Z^*_{5}$ из \examplename{} \ref{example:cyclic_group_F5*}. Мы знаем из \examplename{} \ref{example:factor_groupds_of_F*5}, что элементы $2$ и $3$ являются образующими $\Z^*_{5}$, и следовательно, следующее отображение является хеш-функцией Педерсена:
\begin{equation}
H_{\{2,3\}}: \Z_4 \times \Z_4 \to \Z^*_{5}\;;\; (x,y)\mapsto 2^x \cdot 3^y
\end{equation}

Чтобы увидеть, как может быть вычислено это отображение, выберем входное значение $(1,3)$ из $\Z_4 \times \Z_4$. Тогда, используя таблицу умножения из \eqref{Z5_tables}, мы вычисляем $H_{\{2,3\}}(1,3)= 2^1\cdot 3^3= 2\cdot 2 =4$. 

Чтобы увидеть, как композиция хеш-функции с $H_{\{2,3\}}$ определяет хеш-в-группу функцию, рассмотрим хеш-функцию $SHA256$ из примера \ref{ex:SHA256}. Дана некоторая двоичная строка $s\in\{0,1\}^*$, мы можем вставить два наименее значимых бита $SHA256(s)_0$ и $SHA256(s)_1$ из образа $SHA256(s)$ в $H_{\{2,3\}}$, чтобы получить элемент в $\F_5^*$. Это определяет следующую хеш-в-группу функцию
$$
SHA256\_H_{\{2,3\}}: \{0,1\}^* \to \Z_5^*\;;\; s \mapsto 2^{SHA256(s)_0}\cdot 3^{SHA256(s)_1}
$$
Чтобы увидеть, как может быть вычислена эта хеш-функция, рассмотрим пустую строку $<>$. Поскольку мы знаем из вычисления Sage в \examplename{} \ref{ex:SHA256}, что $SHA256(<>)_0=1$ и что $SHA256(<>)_1=0$, мы получаем $SHA\_256H_{\{2,3\}}(<>)= 2^1 \cdot 3^0 = 2$. 

Конечно, вычисление $SHA256\_H_{\{2,3\}}$ в стиле <<бумага и ручка>> затруднительно. Однако мы можем легко реализовать эту функцию в Sage следующим образом:
\begin{sagecommandline}
sage: import hashlib
sage: def SHA256_H(x):
....:     Z5 = Integers(5) # определяем групповой тип
....:     hasher = hashlib.sha256(x) # вычисляем SHA256
....:     digest = hasher.hexdigest()
....:     z = ZZ(digest, 16) # приводим к целому
....:     z_bin = z.digits(base=2, padto=256) # приводим к 256 битам
....:     return Z5(2)^z_bin[0] * Z5(3)^z_bin[1]
sage: SHA256_H(b"") # вычисляем на пустой строке
sage: SHA256_H(b"SHA") # возможные образы -- {1,2,3}
sage: SHA256_H(b"Math")
\end{sagecommandline}
\end{example}
\begin{exercise}
\label{exercise:Pedersen_hash_1}
Рассмотрим мультипликативную группу $\Z_{13}^*$ арифметики по модулю $13$ из \examplename{} \ref{ex:Zn*}. Выберите множество из $3$ образующих $\Z_{13}^*$, определите соответствующую \capitalisewords{хеш-функцию Педерсена} и вычислите \capitalisewords{хеш Педерсена} для $(3,7,11)\in \Z_{12}$.
\end{exercise}
\begin{exercise}
Рассмотрите \capitalisewords{хеш Педерсена} из \exercisename{} \ref{exercise:Pedersen_hash_1}. Скомпонуйте его с хеш-функцией $SHA256$ из \examplename{} \ref{ex:SHA256}, чтобы определить хеш-в-группу функцию. Реализуйте эту функцию в Sage.
\end{exercise}

\subsubsection{Семейства псевдослучайных функций в DDH-безопасных группах}
Как отмечено в \ref{def:Pedersen_hash}, семейство хеш-функций Педерсена, параметризованное множеством образующих $\{g_1,\ldots,g_j\}$, не квалифицируется как семейство псевдослучайных функций, и поэтому не должно быть инстанцировано как таковое. Чтобы увидеть пример правильного семейства псевдослучайных функций в группах, где предположение диффи -- хеллмана о различении (см. \secname \ref{def:DDH-secure}) предполагается верным, пусть $\G$ -- DDH-безопасная циклическая группа порядка $n$ с образующим $g$, и пусть $\{a_0,a_1,\ldots,a_k\}\subset \Z_{n}^*$ -- равномерно случайно сгенерированное множество чисел, обратимых в арифметике по модулю $n$. Тогда семейство псевдослучайных функций, параметризованное \uterm{зерном} (seed) $\{a_0,a_1,\ldots,a_k\}$, задаётся следующим образом:
\begin{equation}
\label{prf_in_cyclic_group}
F_{\{a_0,a_1,\ldots,a_k\}}: \{0,1\}^{k+1} \to \G:\; (b_0,\ldots,b_k)\mapsto g^{\Pi_{i=0}^k a_i^{b_i}}
\end{equation}
\begin{exercise} Рассмотрим мультипликативную группу $\Z_{13}^*$ арифметики по модулю $13$ из \examplename{} \ref{ex:Zn*} и параметр $k=2$. Выберите образующий $\Z_{13}^*$, зерно и \uterm{инстанцируйте} (instantiate) член семейства, заданного в \eqref{prf_in_cyclic_group}, для этого зерна. Вычислите этот член на двоичной строке $<1,0,1>$.
\end{exercise}

\section{Коммутативные кольца}\label{sec:rings}
В предыдущем разделе мы видели, что целые числа являются коммутативной группой относительно сложения целых чисел. Однако, как нам известно, на целых числах определены две арифметические операции: сложение и умножение. Однако, в отличие от сложения, умножение не определяет групповую структуру, поскольку целые числа обычно не имеют мультипликативных обратных. Конфигурации такого рода составляют так называемые \term{коммутативные кольца с единицей} (commutative rings with unit), и определяются следующим образом:

\begin{definition}[Коммутативное кольцо с единицей]\label{def:comm-ring-unit}
\term{Коммутативное кольцо с единицей} $ (R, +, \cdot, 1) $ (commutative ring with unit) -- это множество $R$ с двумя отображениями, $ +: R \times R \to R $ и $ \cdot: R \times R \to R $, называемыми \term{сложением} (addition) и \term{умножением} (multiplication), и элементом $1\in R$, называемым \term{единицей} (unit), такими, что выполняются следующие условия:
\begin{itemize}
\item $ (R, +) $ -- коммутативная группа, где нейтральный элемент обозначается $ 0 $.
\item \hilight{Коммутативность умножения} (Commutativity of multiplication): $r_1\cdot r_2 = r_2\cdot r_1$ для всех $r_1, r_2\in R$.
\item \hilight{Мультипликативная нейтральная единица} (Multiplicative neutral unit): $1\cdot g=g$ для всех $g\in R$.
\item \hilight{Ассоциативность} (Associativity): Для всех $g_1,g_2,g_3\in\R$ выполняется равенство
$g_1\cdot(g_2\cdot g_3) = (g_1\cdot g_2)\cdot g_3$.
\item \hilight{Дистрибутивность} (Distributivity): Для всех $ g_1, g_2, g_3 \in\R $ выполняется дистрибутивный закон
$ g_1 \cdot \left (g_2 + g_3 \right) = g_1 \cdot g_2 + g_1 \cdot g_3$.
\end{itemize}
Если $(R,+,\cdot,1)$ -- коммутативное кольцо с единицей, и $R'\subseteq R$ -- подмножество $R$ такое, что ограничение сложения и умножения на $R'$ определяет коммутативное кольцо со сложением $+: R'\times R' \to R'$, умножением $\cdot: R'\times R' \to R'$ и единицей $1$ на $R'$, то $(R',+,\cdot,1)$ называется \term{подкольцом} (subring) кольца $(R,+,\cdot,1)$.
\end{definition}
\begin{notation}Поскольку в этой книге мы рассматриваем исключительно коммутативные кольца, мы часто просто называем их кольцами, оставляя коммутативность неявной.
Множество $R$ с двумя отображениями, $+$ и $\cdot$, которое удовлетворяет всем перечисленным выше правилам, кроме закона коммутативности умножения, называется некоммутативным кольцом. 

Если нет риска неоднозначности (относительно того, каковы отображения сложения и умножения кольца), мы часто опускаем символы $+$ и $\cdot$ и просто пишем $R$ как обозначение кольца, оставляя эти отображения неявными. В этом случае мы также говорим, что $R$ имеет кольцевой тип, указывая на то, что $R$ -- это не просто множество, а множество вместе с отображениями сложения и умножения.\footnote{Коммутативные кольца -- обширная область исследований в математике, и существует бесчисленное множество книг по этой теме. Для наших целей введение дано в \chaptname{} 1, \secname{} 2 книги \cite{nieder-1986}.}
\end{notation}
\begin{example}[Кольцо целых чисел] Множество $\Z$ целых чисел с обычным сложением и умножением -- это архетипический пример коммутативного кольца с единицей $1$. 
\begin{sagecommandline}
sage: ZZ
\end{sagecommandline}
\end{example}
\begin{example}[Лежащая в основе коммутативная группа кольца] Каждое коммутативное кольцо с единицей $(R,+,\cdot,1)$ порождает группу, если мы игнорируем умножение.
\end{example}
Следующий пример несколько необычен, но мы призываем вас продумать его, поскольку это помогает отвлечься от привычных стилей вычислений и сосредоточиться на абстрактном алгебраическом объяснении.
\begin{example} Пусть $S:=\{\bullet,\star,\odot,\otimes\}$ -- множество, содержащее четыре элемента, и пусть сложение и умножение на $S$ определены следующим образом:
\begin{equation}\label{moon-table}
  \begin{tabular}{c | c c c c c c}
    $\cup$ & $\bullet$ & $\star$ & $\odot$ & $\otimes$ \\\\\hline
    $\bullet$ & $\bullet$ & $\star$ & $\odot$ & $\otimes$ \\
    $\star$ & $\star$ & $\odot$ & $\otimes$ & $\bullet$ \\
    $\odot$ & $\odot$ & $\otimes$ & $\bullet$ & $\star$ \\
    $\otimes$ & $\otimes$ & $\bullet$ & $\star$ & $\odot$ \\
  \end{tabular} \quad \quad \quad \quad
  \begin{tabular}{c | c c c c c c}
$ \circ $ & $\bullet$ & $\star$ & $\odot$ & $\otimes$ & \\\\\hline
        $\bullet$ & $\bullet$ & $\bullet$ & $\bullet$ & $\bullet$ &\\
        $\star$ & $\bullet$ & $\star$ & $\odot$ & $\otimes$ &\\
        $\odot$ & $\bullet$ & $\odot$ & $\bullet$ & $\odot$ &\\
        $\otimes$ & $\bullet$ & $\otimes$ & $\odot$ & $\star$ &\\
  \end{tabular}
\end{equation}
Тогда $(S,\cup,\circ, \star)$ -- кольцо с единицей $\star$ и нулём $\bullet$. Поэтому имеет смысл искать решения уравнений, подобных следующему:
\begin{equation}\label{eq:moon}
\otimes \circ (x \cup \odot ) = \star
\end{equation}

Задача здесь состоит в том, чтобы найти $x\in S$ такое, что выполняется \eqref{eq:moon}. Чтобы понять, как такое <<moonmath-уравнение>> может быть решено, мы должны помнить, что кольца ведут себя в основном как обычные числа, когда дело доходит до скобок и правил вычислений. Единственные различия -- это символы и фактический способ сложения и умножения. Имея это в виду, мы решаем уравнение для $x$ <<обычным способом>>: \footnote{Отметим, что существуют более эффективные способы решения этого уравнения. Цель нашего вычисления -- показать, как аксиомы кольца могут быть использованы для решения уравнения.}
\begin{align*}
\otimes \circ (x \cup \odot ) &= \star & \text{ \# применяем дистрибутивный закон}\\
\otimes \circ x \cup \otimes \circ \odot  &= \star &\# \otimes \circ \odot = \odot\\
\otimes \circ x \cup \odot  &= \star & \text{\# прибавляем обратный по $\cup$ элемент $\odot$ к обеим сторонам}\\
\otimes \circ x \cup \odot \cup -\odot  &= \star \cup -\odot & \# \odot \cup -\odot = \bullet\\
\otimes \circ x \cup \bullet &= \star \cup -\odot & \text{\# $\bullet$ -- нейтральный элемент по $\cup$}\\
\otimes \circ x &= \star \cup -\odot & \text{\# для $\cup$ имеем $-\odot = \odot$} \\
\otimes \circ x &= \star \cup \odot &\# \star \cup \odot = \otimes \\
\otimes \circ x &= \otimes  &\text{\# прибавляем обратный по $\circ$ элемент $\otimes$ к обеим сторонам}\\
(\otimes)^{-1}\circ \otimes \circ x &= (\otimes)^{-1}\circ \otimes & \text{\# умножаем на мультипликативный обратный}\\
\star \circ x &= \star\\
x &= \star
\end{align*}
Несмотря на то что это уравнение на первый взгляд выглядело совершенно чужеродным, мы смогли решить его по сути точно так же, как решаем <<обычные>> уравнения, содержащие числа.

Отметим, однако, что всякий раз, когда для решения уравнения обычным способом в кольце требуется мультипликативный обратный, всё может быть совсем не так, как мы привыкли. Например, следующее уравнение нельзя решить для $x$ обычным способом, поскольку в нашем кольце нет мультипликативного обратного для $\odot$.

\begin{equation}
\odot \circ x = \otimes
\end{equation}

Мы можем подтвердить это, посмотрев на таблицу умножения в \eqref{moon-table}, чтобы увидеть, что такого $x$ не существует.

В качестве другого примера, следующее уравнение имеет не одно решение, а два: $x\in\{\star, \otimes\}$.

\begin{equation}
\odot \circ x = \odot
\end{equation}

Отсутствие решения или наличие двух решений -- это определённо не то, к чему мы привыкли от типов вроде рациональных чисел $\mathbb{Q}$.
\end{example}

\begin{example}[Кольцо многочленов] Рассмотрев снова определение многочленов из раздела \ref{sec:polynomial_arithmetics}, мы замечаем, что то, что мы неформально назвали типом $R$ коэффициентов, на самом деле должен быть коммутативным кольцом с единицей, поскольку нам нужны сложение, умножение, коммутативность и существование единицы, чтобы $R[x]$ имел ожидаемые свойства.

На самом деле, если мы считаем $R$ кольцом и определяем сложение и умножение многочленов, как в \eqref{def:polynomial_arithmetic}, множество $R[x]$ является коммутативным кольцом с единицей, где многочлен $1$ является мультипликативной единицей. Мы называем это кольцо \term{кольцом многочленов с коэффициентами в} $R$ (ring of polynomials with coefficients in $R$).
\begin{sagecommandline}
sage: ZZ['x']
\end{sagecommandline}
\end{example}
\begin{example}[Кольцо модульной арифметики $n$]
\label{def:ring_of_mod_n_arithmetics}
 Пусть $n$ -- модуль и $(\Z_n,+,\cdot)$ -- множество всех классов остатков целых чисел по модулю $n$ с проекцией сложения и умножения целых чисел, как определено в \secname \ref{def:remainder_class_representation}. Тогда $(\Z_n,+,\cdot)$ -- коммутативное кольцо с единицей $1$.
\begin{sagecommandline}
sage: Integers(6)
\end{sagecommandline}
\end{example}
\begin{example}[Двоичные представления в модульной арифметике]
\smecomb{TODO}{add example} (Нет единственности)
\end{example}

\begin{example}[Вычисление многочлена в показателе степени образующих групп] Как мы показываем в \secname{} \ref{sec:QAP}, ключевой идеей во многих протоколах с нулевым разглашением является возможность кодирования вычислений как многочленов, а затем сокрытие информации о том вычислении путём вычисления многочлена <<в показателе степени>> определённых криптографических групп (\secname{} \ref{sec:gorth_16}).

Чтобы понять основной принцип этой идеи, рассмотрим снова экспоненциальное отображение \eqref{exponentialmap}. Если $\G$ -- конечная циклическая группа порядка $n$ с образующим $g\in\G$, то кольцевая структура $(\Z_n,+,\cdot)$ соответствует групповой структуре $\G$ следующим образом:
\begin{align}
\label{def:ring_exponential_laws}
g^{x+y} &= g^x\cdot g^y &
g^{x\cdot y} &= \left( g^x\right)^y & \text{для всех } x,y\in\Z_n
\end{align}
Это соответствие позволяет вычислять многочлены с коэффициентами в $\Z_n$ <<в показателе степени>> образующего группы. Чтобы понять, что это означает, пусть $p\in \Z_n[x]$ -- многочлен с $p(x)=a_m\cdot x^m+a_{m-1}x^{m-1}+\ldots + a_1x +a_0$, и пусть $s\in\Z_n$ -- точка вычисления. Тогда ранее определённые экспоненциальные законы \ref{def:ring_exponential_laws} влекут следующее тождество:
\begin{align}
\label{def:polynomial_ring_exponential_laws}
g^{p(s)} & = g^{a_m\cdot s^m+a_{m-1}s^{m-1}+\ldots + a_1s +a_0}\
         & = \left(g^{s^m}\right)^{a_m}\cdot \left(g^{s^{m-1}}\right)^{a_{m-1}}\cdot \ldots\cdot \left(g^{s}\right)^{a_1}\cdot g^{a_0}\notag
\end{align}
Используя эти тождества, возможно вычислить любой многочлен $p$ степени $deg(p)\leq m$ в <<секретной>> точке вычисления $s$ в показателе степени $g$ без каких-либо знаний о $s$, при условии, что $\G$ является DL-группой. Чтобы увидеть это, предположим, что дано множество $\{g,g^s, g^{s^2},\ldots, g^{s^m}\}$, но $s$ неизвестно. Тогда
$g^{p(s)}$ может быть вычислено с использованием \eqref{def:polynomial_ring_exponential_laws}, но вычислить $s$ невозможно.
\end{example}

\begin{example} Чтобы привести пример вычисления многочлена в показателе степени конечной циклической группы, рассмотрим экспоненциальное отображение из \examplename{} \ref{ex:in-the-exponent}:

\begin{equation}
3^{(\cdot)}: \Z_4 \to \Z_5^* \;;\; x \mapsto 3^x
\end{equation}

Выбирая многочлен $p(x)= 2x^2 +3x +1$ из $\Z_4[x]$, мы сначала вычисляем многочлен в точке $s=2$, а затем записываем результат в показатель степени $3$ следующим образом:
\begin{align*}
3^{p(2)} &=3^{2\cdot 2^2+3\cdot 2 +1}\
          & = 3^{2\cdot 0 +2 +1}\
          & = 3^{3}\
          & = 2
\end{align*}
Это было возможно, потому что у нас был доступ к точке вычисления $2$. С другой стороны, если бы у нас был доступ только к множеству $\{3, 4, 1\}$, и мы знали, что это множество представляет множество $\{3,3^s, 3^{s^2}\}$ для некоторого секретного значения $s$, мы могли бы вычислить
$p$ в точке $s$ в показателе степени $3$ следующим образом:
\begin{align*}
3^{p(s)} &= 1^2 \cdot 4^3\cdot 3^1\
         &= 1\cdot 4\cdot 3\
         &= 2
\end{align*}
Оба вычисления \comms{согласуются}, поскольку секретная точка $s$ была равна $2$ в этом примере. Однако второе вычисление было возможно без каких-либо знаний о $s$.
\end{example}

\subsection{Хеширование в модульную арифметику}
\label{hash-to-modular-arithmetics}
Как мы видели в \secname{} \ref{sec:hashing-to-groups}, известны и используются в приложениях различные конструкции хеширования в группы. Поскольку коммутативные кольца являются коммутативными группами, когда мы игнорируем умножение, конструкции хеш-в-группу могут применяться для хеширования в коммутативные кольца. Мы рассмотрим некоторые часто используемые приложения ниже.

\sme{put subsubsection title here}

Одним из наиболее широко используемых приложений конструкций хеш-в-кольцо являются хеш-функции, отображающие в кольцо $\Z_n$ арифметики по модулю $n$ для некоторого модуля $n$. Известны различные подходы к построению такой функции, но, вероятно, наиболее широко используемые основаны на интуиции, что образы общих хеш-функций могут быть интерпретированы как двоичные представления целых чисел, как объяснено в \examplename{} \ref{naive-cyclic-group-hash}.

Из этой интерпретации следует, что один простой метод хеширования в $\Z_n$ строится на наблюдении, что если $n$ -- модуль с битовой длиной \eqref{def:binary_representation_integer} $k=|n|$, то каждая двоичная строка $<b_0,b_1,\ldots,b_{k-2}>$ длины $k-1$ определяет целое число $z$ в диапазоне $0\leq z \leq 2^{k-1}-1< n $:
\begin{equation}
z = b_0\cdot 2^0 + b_1\cdot 2^1 + \ldots + b_{k-2}\cdot 2^{k-2}
\end{equation}
Теперь, поскольку $z<n$, мы знаем, что $z$ гарантированно находится в множестве $\{0,1,\ldots,n-1\}$, и следовательно, оно может быть интерпретировано как элемент $\Z_n$. Следовательно, если $H:\{0,1\}^*\to\{0,1\}^{k-1}$ -- хеш-функция, то хеш-в-кольцо функция может быть построена следующим образом:
\begin{equation}\label{eq:hash-Zr}
H_{|n|_2-1}: \{0,1\}^* \to \Z_r: \; s \mapsto
H(s)_0\cdot 2^0 + H(s)_1\cdot 2^1 + \ldots + H(s)_{k-2}\cdot 2^{k-2}
\end{equation}

Недостаток этой хеш-функции заключается в том, что распределение хеш-значений в $\Z_n$ не обязательно равномерно. На самом деле, если $n$ больше $2^{k-1}$, то по конструкции $H_{|n|_2-1}$ никогда не будет хешировать на значения $z\geq 2^{k-1}$. Поэтому использование этого метода хеширования генерирует примерно равномерные хеши только если $n$ очень близко к $2^{k-1}$. В худшем случае, когда $n=2^k-1$, она упускает почти половину всех элементов из $\Z_n$.

Преимуществом этого подхода является то, что свойства, такие как стойкость к нахождению прообраза или стойкость к коллизиям (см. \secname{} \ref{sec:hash-functions}) исходной хеш-функции $H(\cdot)$, сохраняются.
\begin{example} Чтобы исследовать равномерность хеширования в $\Z_n$ с использованием метода, описанного в \ref{eq:hash-Zr}, рассмотрим модуль $n$, представимый как 5-битное двоичное число, что указывает на то, что $n$ -- целое число в диапазоне $16 \leq n < 32$.

Наиболее равномерное распределение хешей достигается, когда $n = 16$, потому что кольцо $\Z_{16}$ состоит из элементов $\{0, 1, \ldots, 15\}$. В этом сценарии мы можем использовать хеш-функцию $H_{|n|2-1}$, усечённую до первых $4$ битов $SHA256$, как продемонстрировано в \examplename{} \ref{ex:SHA256}. Это позволяет нам определить хеш-функцию в $\Z_{16}$ следующим образом:
$$
H_{|16|_2-5}: \{0,1\}^* \to \Z_{16}:\; s\mapsto
SHA256(s)_0\cdot 2^0 + SHA256(s)_1\cdot 2^1 + \ldots + SHA256(s)_3\cdot 2^3
$$
Поскольку $k=|16|_2=5$ и $16-2^{k-1}=0$, этот хеш отображает равномерно на $\Z_{16}$. Мы можем использовать Sage для её реализации:
\begin{sagecommandline}
sage: import hashlib
sage: def Hash5(x):
....:     hasher = hashlib.sha256(x) # вычисляем SHA256
....:     digest = hasher.hexdigest()
....:     d = ZZ(digest, base=16) # приводим к целому
....:     d = d.str(2)[-4:] # оставляем 4 наименее значимых бита
....:     d = ZZ(d, base=2) # приводим к целому
....:     return d
sage: Hash5(b'')
\end{sagecommandline}
Мы можем затем использовать Sage, чтобы применить эту функцию к большому набору входных значений для построения визуализации распределения по множеству $\{0,\ldots,15\}$. Выполнение на более чем $500$ входных значений даёт следующий график:
\begin{sagesilent}
H1 = list_plot([Hash5(ZZ(k).str(2).encode('utf-8')) for k in range(500)])
\end{sagesilent}
\begin{center}
\sageplot[scale=.5]{H1}
\end{center}
Чтобы получить интуицию о равномерности, мы можем подсчитать число раз, когда хеш-функция $H_{|16|_2-1}$ отображает на каждое число в множестве $\{0,1,\ldots,15\}$ в цикле из $100000$ хешей, и сравнить это с идеальным равномерным распределением, которое отображало бы ровно 6250 выборок на каждый элемент. Это даёт следующий результат:
\begin{sagesilent}
arr = []
arr = [0 for i in range(16)]
for i in range(100000):
    arr[Hash5(ZZ(i).str(2).encode('utf-8'))] +=1
H2 = list_plot(arr, ymin=0,ymax=10000)
\end{sagesilent}
\begin{center}
\sageplot[scale=.5]{H2}
\end{center}
Неравномерность становится очевидной, если мы хотим построить подобную хеш-функцию для $\Z_n$ для любого другого 5-битного целого числа $n$ в диапазоне $17\leq n < 32$. В этом случае определение хеш-функции точно такое же, как для $\Z_{16}$, и следовательно, образы не будут превышать значение $15$. Так, например, даже в случае хеширования в $\Z_{31}$ хеш-функция никогда не отображает ни на какое значение больше $15$, оставляя почти половину всех чисел вне диапазона образа.
\begin{sagesilent}
arr = []
arr = [0 for i in range(31)]
for i in range(100000):
    arr[Hash5(ZZ(i).str(2).encode('utf-8'))] +=1
H3 = list_plot(arr, ymin=0,ymax=10000)
\end{sagesilent}
\begin{center}
\sageplot[scale=.5]{H3}
\end{center}
\end{example}

\sme{put subsubsection title here}

Второй широко используемый метод хеширования в $\Z_n$ строится на следующем наблюдении: Если $n$ -- модуль с битовой длиной $|n|_2=k_1$, а $H:\{0,1\}^*\to \{0,1\}^{k_2}$ -- хеш-функция, производящая дайджесты размера $k_2$, и $k_2\geq k_1$, то хеш-в-кольцо функция может быть построена интерпретацией образа $H$ как двоичного представления целого числа, а затем взятием модуля по $n$ для отображения в $\Z_n$:
\begin{equation}
H'_{mod_n}: \{0,1\}^* \to \Z_n: \; s \mapsto
\Zmod{\left(H(s)_0\cdot 2^0 + H(s)_1\cdot 2^1 + \ldots + H(s)_{k_2}\cdot 2^{k_2}\right)}{n}
\end{equation}

Недостаток этой хеш-функции заключается в том, что вычисление модуля требует некоторых вычислительных затрат. Кроме того, распределение хеш-значений в $\Z_n$ может быть неравномерным, в зависимости от числа $\Zmod{2^{k_2+1}}{n}$. Преимуществом этой функции является то, что потенциальные свойства исходной хеш-функции $H(\cdot)$ (такие как стойкость к нахождению прообраза или стойкость к коллизиям) сохраняются, и распределение может быть сделано почти равномерным, с лишь пренебрежимо малым смещением в зависимости от того, какие модуль $n$ и размер образа $k_2$ выбраны.
\begin{example} Чтобы дать реализацию хеш-функции $H_{mod_n}$, мы используем усечение $SHA256$ до $k_2$ битов из \examplename{} \ref{ex:SHA256}, и определяем хеш в $\Z_{23}$ следующим образом:
\begin{multline*}
H_{mod_{23},k_2}: \{0,1\}^* \to \Z_{23}:\; \\
s\mapsto
\Zmod{\left(SHA256(s)_0\cdot 2^0 + SHAH256(s)_1\cdot 2^1 + \ldots + SHA256(s)_{k_2}\cdot 2^{k_2}\right)}{23}
\end{multline*}
Мы хотим использовать различные инстанциации $k_2$ для визуализации влияния длины усечения на распределение хешей в $\Z_{23}$. Мы можем использовать Sage для её реализации следующим образом:
\begin{sagecommandline}
sage: import hashlib
sage: Z23 = Integers(23)
sage: def Hash_mod23(x, k2):
....:     hasher = hashlib.sha256(x.encode('utf-8')) # Вычисляем SHA256
....:     digest = hasher.hexdigest()
....:     d = ZZ(digest, base=16) # приводим к целому
....:     d = d.str(2)[-k2:] # оставляем k2+1 наименее значимых битов
....:     d = ZZ(d, base=2) # приводим к целому
....:     return Z23(d) # приводим к Z23
\end{sagecommandline}

Мы можем затем использовать Sage, чтобы применить эту функцию к большому набору входных значений для построения визуализаций распределения по множеству $\{0,\ldots,22\}$ для различных значений $k_2$, подсчитывая число раз, когда она отображает на каждое число в цикле из $100000$ хешей. Мы получаем следующий график:
\begin{sagesilent}
arr1 = []
arr1 = [0 for i in range(23)]
for i in range(100000):
    arr1[Hash_mod23(ZZ(i).str(2),5)] +=1
H3 = list_plot(arr1, ymin=0,ymax=10000,color='red', legend_label='k2=5')
arr2 = []
arr2 = [0 for i in range(23)]
for i in range(100000):
    arr2[Hash_mod23(ZZ(i).str(2),7)] +=1
H4 = list_plot(arr2, ymin=0,ymax=10000,color='blue', legend_label='k2=7')
arr3 = []
arr3 = [0 for i in range(23)]
for i in range(100000):
    arr3[Hash_mod23(ZZ(i).str(2),9)] +=1
H5 = list_plot(arr3, ymin=0,ymax=10000,color='yellow', legend_label='k2=9')
arr4 = []
arr4 = [0 for i in range(23)]
for i in range(100000):
    arr4[Hash_mod23(ZZ(i).str(2),16)] +=1
H6 = list_plot(arr4, ymin=0,ymax=10000,color='black', legend_label='k2=16')
\end{sagesilent}
\begin{center}
\sageplot[scale=.6]{H3+H4+H5+H6}
\end{center}
\end{example}

\subsubsection{Метод <<попробуй и приращивай>>}\label{def:try_and_increment_hash}

Третий метод, который иногда можно встретить в реализациях, -- это так называемый \term{метод <<попробуй и приращивай>>} (``try-and-increment'' method). Чтобы понять этот метод, мы определяем целое число $z\in\Z$ из любого хеш-значения $H(s)$, как мы делали в предыдущих методах:

\begin{equation}
z = H(s)_0\cdot 2^0 + H(s)_1\cdot 2^1 + \ldots + H(s)_{k}\cdot 2^{k}
\end{equation}

Хеширование в $\Z_n$ затем достижимо сначала вычислением $z$, а затем проверкой, принадлежит ли $z\in\Z_n$. Если да, то хеширование выполнено; если нет, строка $s$ модифицируется детерминированным образом, и процесс повторяется, пока не будет найден подходящий элемент $z\in\Z_n$. Подходящей детерминированной модификацией может быть конкатенация исходной строки с некоторым битовым счётчиком. Алгоритм <<попробуй и приращивай>> тогда работает, как в \algname{} \ref{alg_try_and_increment}.
\begin{algorithm}\caption{Хеш-в-$\Z_n$}
\label{alg_try_and_increment}
\begin{algorithmic}[0]
\Require $n \in \Z$ с $|n|_2=k$ и $s\in\{0,1\}^*$
\Procedure{Try-and-Increment}{$n,k,s$}
\State $c \gets 0$
\Repeat
\State $s' \gets s||c\_bits()$
\State $z \gets H(s')_0\cdot 2^0 + H(s')_1\cdot 2^1 + \ldots + H(s')_{k}\cdot 2^{k}$
\State $c\gets c+1$
\Until{$z<n$}
\State \textbf{return} $x$
\EndProcedure
\Ensure $ z\in \Z_n$
\end{algorithmic}
\end{algorithm}

В зависимости от параметров этот метод может быть очень эффективным. На самом деле, если $k$ достаточно велико и $n$ близко к $2^{k+1}$, вероятность $z<n$ очень высока, и цикл повторения почти всегда будет выполнен только один раз. Недостатком, однако, является то, что вероятность необходимости выполнить цикл несколько раз не равна нулю.

\section{Поля}\label{sec:fields}
Мы начали эту главу с определения группы (\secname{} \ref{sec:groups}), которое затем расширили до определения коммутативного кольца с единицей (\secname \ref{sec:rings}). Эти типы колец обобщают поведение целых чисел. В этом разделе мы рассмотрим те частные случаи коммутативных колец, где каждый элемент, отличный от нейтрального элемента сложения, имеет мультипликативный обратный. Эти структуры ведут себя очень похоже на множество рациональных чисел $\mathbb{Q}$. Рациональные числа являются, в некотором смысле, расширением кольца целых чисел, то есть они построены включением вновь определённых мультипликативных обратных (дробей) к целым числам. Поля определяются следующим образом:

\begin{definition}[Поле]\label{def:field}
\term{Поле} $ (\F, +, \cdot) $ (field) -- это множество $\F$ с двумя отображениями $ +: \F \times \F \to \F $ и $ \cdot: \F \times \F \to \F $, называемыми \term{сложением} (addition) и \term{умножением} (multiplication), такими, что выполняются следующие условия:
\begin{itemize}
\item $ \left (\F, + \right) $ -- коммутативная группа, где нейтральный элемент обозначается $ 0 $.
\item $ \left (\F \setminus \left \{0 \right \}, \cdot \right) $ -- коммутативная группа, где нейтральный элемент обозначается $ 1 $.
\item (Дистрибутивность) Уравнение $g_1 \cdot \left (g_2 + g_3 \right) = g_1 \cdot g_2 + g_1 \cdot g_3$ выполняется для всех $ g_1, g_2, g_3 \in \F $.
\end{itemize}
Если $(\F,+,\cdot)$ -- поле и $\F'\subset \F$ -- подмножество $\F$ такое, что ограничение сложения и умножения на $\F'$ определяет поле со сложением $+: \F'\times \F' \to \F'$ и умножением $\cdot: \F'\times \F' \to \F'$ на $\F'$, то $(\F',+,\cdot)$ называется \term{подполем} (subfield) поля $(\F,+,\cdot)$, а $(\F,+,\cdot)$ называется \term{расширением поля} (extension field) поля $(\F',+,\cdot)$.
\end{definition}

\begin{notation} Если нет риска неоднозначности (относительно того, каковы отображения сложения и умножения поля), мы часто опускаем символы $+$ и $\cdot$ и просто пишем $\F$ как обозначение поля, оставляя отображения неявными. В этом случае мы также говорим, что $\F$ имеет полевой тип, указывая на то, что $\F$ -- это не просто множество, а множество со сложением и умножением, удовлетворяющее определению поля (\ref {def:field}).\footnote{Поскольку поля имеют большое значение в криптографии и теории чисел, существует множество книг по этой теме. Для общего введения см., например, \chaptname{} 6, \secname{} 1 в \cite{mignotte-1992}, или \chaptname{} 1, \secname{} 2 в \cite{nieder-1986}.}

Мы называем $(\F,+)$ \term{аддитивной группой} (additive group) поля. Мы используем обозначение $\F^*:= \F \setminus \left \{0 \right \}$ для множества всех элементов, исключая нейтральный элемент сложения, и называем $(\F^*,\cdot)$ \term{мультипликативной группой} (multiplicative group) поля.
\end{notation}

\term{Характеристикой}\label{def:characteristic} поля $ \F $, обозначаемой $char(\F)$, является наименьшее натуральное число $ n \geq 1 $, для которого $n$-кратная сумма мультипликативного нейтрального элемента $ 1 $ равна нулю, т.е. для которого $ \sum_{i = 1} ^ n 1 = 0 $. Если такое $ n> 0 $ существует, поле называется имеющим \term{конечную характеристику} (finite characteristic). Если, с другой стороны, каждая конечная сумма $1$ такова, что она не равна нулю, то поле по определению имеет характеристику $ 0 $.\sme{Check change of wording} \smelong{S: Tried to disambiguate the scope of negation between 1. ``It is true of every finite sum of $1$ that it is not equal to 0''  and 2. ``It is not true of every finite sum of $1$ that it is  equal to 0'' From the example below, it looks like 1. is the intended meaning here, correct?}
\begin{example}[Поле рациональных чисел] Вероятно, наиболее известным примером поля является множество рациональных чисел $\mathbb{Q}$ вместе с обычным определением сложения, вычитания, умножения и деления. Поскольку не существует натурального числа $n\in \N$ такого, что $\sum_{j=0}^n 1 =0$ в множестве рациональных чисел, характеристика поля $\mathbb{Q}$ равна $char(\mathbb{Q})=0$.
\begin{sagecommandline}
sage: QQ
\end{sagecommandline}
\end{example}
\begin{example}[Поле с двумя элементами]\label{ex:field-2-elements} Можно показать, что в любом поле нейтральный элемент сложения $0$ должен отличаться от нейтрального элемента умножения $1$, то есть $0\neq 1$ всегда выполняется в поле. Это означает, что наименьшее поле должно содержать по крайней мере два элемента. Как показывают следующие таблицы сложения и умножения, действительно существует поле с двумя элементами, которое обычно называется $\F_2$:

Пусть $\F_2:=\{0,1 \}$ -- множество, содержащее два элемента, и пусть сложение и умножение на $\F_2$ определены следующим образом:
\begin{equation}
  \begin{tabular}{c | c c}
    + & 0 & 1 \\\\\hline
    0 & 0 & 1 \\
    1 & 1 & 0 \\
  \end{tabular} \quad \quad \quad \quad
  \begin{tabular}{c | c c}
$\cdot$ & 0 & 1 \\\\\hline
      0 & 0 & 0 \\
      1 & 0 & 1 \\
  \end{tabular}
\end{equation}
Поскольку $1+1=0$ в поле $\F_2$, мы знаем, что характеристика $\F_2$ равна $char(\F_2)=2$.  Мультипликативная подгруппа $\F_2^*$ поля $\F_2$ задаётся тривиальной группой $\{1\}$.
\begin{sagecommandline}
sage: F2 = GF(2)
sage: F2(1) # Получаем элемент из GF(2)
sage: F2(1) + F2(1) # Сложение
sage: F2(1) / F2(1) # Деление
\end{sagecommandline}
\end{example}
\begin{exercise}
Рассмотрим кольцо арифметики по модулю $5$ $(\Z_5,+,\cdot)$ из \examplename{} \ref{primfield_z_5}. Покажите, что $(\Z_5,+,\cdot)$ является полем. Какова характеристика $\Z_5$? Докажите, что уравнение $a\cdot x = b$ имеет только единственное решение $x\in \Z_5$ для любых данных $a,b\in \Z_5^*$.
\end{exercise}
\begin{exercise}
Рассмотрим кольцо арифметики по модулю $6$ $(\Z_6,+,\cdot)$ из \examplename{} \ref{def_residue_ring_z_6}. Покажите, что $(\Z_6,+,\cdot)$ не является полем.
\end{exercise}

\subsection{Простые поля}
\label{prime_fields}
Как мы видели во многих примерах предыдущих разделов, модульная арифметика во многом ведёт себя аналогично обычной арифметике целых чисел. Это связано с тем, что множества классов остатков $\Z_n$ являются коммутативными кольцами с единицами (см. \examplename{} \ref{def:ring_of_mod_n_arithmetics}).

Однако мы также видели в \examplename{} \ref{ex:modulus-prime-group}, что всякий раз, когда модуль является простым числом, каждый класс остатков, отличный от нулевого класса, имеет модульный мультипликативный обратный. Это важное наблюдение, поскольку оно немедленно влечёт, что в случае, когда модуль является простым числом, множество классов остатков $\Z_n$ является не просто кольцом, а фактически \term{полем} (field). Кроме того, поскольку $\sum_{j=0}^n 1 = 0$ в $\Z_n$, мы знаем, что эти поля имеют конечную характеристику $n$.

\begin{notation}[Простые поля]
\label{def:prime_fields}
Пусть $p \in \Prim$ -- простое число и $(\Z_p,+,\cdot)$ -- кольцо арифметики по модулю $p$ (см. \examplename{} \ref{def:ring_of_mod_n_arithmetics}). Чтобы отличать простые поля от произвольных колец модульной арифметики, мы пишем $ (\F_p, +, \cdot) $ для кольца арифметики по модулю $p$ и называем его \term{простым полем} (prime field) характеристики $p$.
\end{notation}

Простые поля являются основой многих современных основанных на алгебре криптографических систем, поскольку они обладают рядом желательных свойств. Одно из них заключается в том, что любое простое поле характеристики $p$ содержит ровно $p$ элементов, которые могут быть представлены на компьютере не более чем $log_2(p)$ битами. С другой стороны, поля вроде рациональных чисел требуют потенциально неограниченного количества битов для любого представления с полной точностью.\footnote{Для подробного введения в теорию простых полей см., например, \chaptname{} 2 в \cite{nieder-1986}, или \chaptname{} 6 в \cite{mignotte-1992}.}

Поскольку простые поля являются частными случаями колец модульной арифметики, сложение и умножение могут быть вычислены сначала выполнением обычного сложения и умножения целых чисел, а затем рассмотрением остатка в евклидовом делении на $p$ как результата. Для любого элемента простого поля $x\in \F_p$ его аддитивный обратный (отрицательный) задаётся $-x=\Zmod{p-x}{p}$. Для $x\neq 0$ мультипликативный обратный всегда существует и равен $x^{-1}=x^{p-2}$. Деление затем определяется умножением на мультипликативный обратный, как объяснено в \secname{} \ref{sec:modular_inverses}. Альтернативно, мультипликативный обратный может быть вычислен с использованием расширенного алгоритма Евклида, как объяснено в \eqref{eq_compute_multiplicative_inverse}.

\begin{example}
Наименьшим полем является поле $\F_2$ характеристики $2$, как мы видели в \examplename{}  \ref{ex:field-2-elements}. Это простое поле простого числа $2$.
\end{example}
\begin{example}
Поле $\F_5$ из \examplename{} \ref{primfield_z_5} является простым полем, как определено его таблицей сложения и умножения \eqref{Z5_tables}.
\end{example}

\begin{example}\label{prime-field-F5}
Чтобы обобщить основные аспекты вычислений в простых полях, рассмотрим простое поле $\F_5$ (\examplename{} \ref{primfield_z_5}) и упростим следующее выражение:
\begin{equation}
\left(\frac{2}{3} - 2\right)\cdot 2
\end{equation}
Первое, на что следует обратить внимание, заключается в том, что, поскольку $\F_5$ -- поле, все правила идентичны правилам, которые мы изучили в школе, имея дело с рациональными, вещественными или комплексными числами. Это означает, что мы можем использовать методы вроде раскрытия скобок (дистрибутивность) или сложение как обычно. Для облегчения вычислений мы можем обратиться к таблицам сложения и умножения в \eqref{Z5_tables}.
\begin{align*}
\left(\frac{2}{3} - 2\right)\cdot 2 &=
 \frac{2}{3}\cdot 2 - 2\cdot 2 & \text{\# дистрибутивный закон}\\
 &= \frac{2\cdot 2}{3} - 2\cdot 2 & \Zmod{4}{5}=4 \\
 &= \frac{4}{3} - 4 & \text{\# мультипликативный обратный 3 есть } \Zmod{3^{5-2}}{5}=2\\
 &= 4\cdot 2 - 4 & \text{\# аддитивный обратный 4 есть } 5-4=1\\
 &= 4\cdot 2 +1 & \Zmod{8}{5}=3\\
 &= 3 +1 & \Zmod{4}{5}=4\\
 &= 4
\end{align*}
В этом примере мы вычислили мультипликативный обратный $3$ с использованием тождества
$x^{-1}=x^{p-2}$ в простом поле. Это непрактично для больших простых чисел. Напомним, что другой способ вычисления мультипликативного обратного -- расширенный алгоритм Евклида (см. \ref{eq: erw_Eukl_algo}).  Чтобы освежить память, алгоритм решает уравнение $x^{-1}\cdot 3 + t \cdot 5 =1$ для $x^{-1}$ (хотя $t$ в данном случае не имеет значения). Мы получаем следующее:
\begin{equation}
  \begin{tabular}{c | c c l}
    k & $ r_k $ & $ x^{-1}_k $ & $ t_k $ \\\\\hline
    0 & 3 & 1 & $\cdot$\ \\
    1 & 5 & 0 & $\cdot$ \\
    2 & 3 & 1 & $\cdot$ \\
    3 & 2 &-1 & $\cdot$ \\
    4 & 1 & 2  & $\cdot$ \\
  \end{tabular}
\end{equation}
Таким образом, мультипликативный обратный $3$ в $\Z_5$ равен $2$, и действительно, если мы вычислим произведение $3$ с его мультипликативным обратным $2$, мы получим нейтральный элемент $1$ в $\F_5$.
\end{example}
\begin{exercise}[Простое поле $\F_3$]\label{exercise:F3} Постройте таблицы сложения и умножения простого поля $\F_3$.
\end{exercise}
\begin{exercise}[Простое поле $\F_{13}$]\label{prime_field_F13} Постройте таблицы сложения и умножения простого поля $\F_{13}$.
\end{exercise}
\begin{exercise} Рассмотрим простое поле $\F_{13}$ из \exercisename{} \ref{prime_field_F13}. Найдите множество всех пар $(x,y)\in \F_{13}\times \F_{13}$, удовлетворяющих следующему уравнению:
\begin{equation}
x^2+y^2 = 1 + 7\cdot x^2\cdot y^2
\end{equation}
\end{exercise}
\subsection{Квадратные корни} Как нам известно из арифметики целых чисел, некоторые целые числа, например $4$ или $9$, являются квадратами других целых чисел: например, $4=2^2$ и $9=3^2$. Однако мы также знаем, что не все целые числа являются квадратами других целых чисел: например, не существует целого числа $x\in\Z$ такого, что $x^2=2$. Если целое число $a$ является квадратом другого целого числа $b$, то имеет смысл определить квадратный корень из $a$ как $b$.

В контексте простых полей элемент, являющийся квадратом другого элемента, также называется \term{квадратичным вычетом} (quadratic residue), а элемент, не являющийся квадратом другого элемента, называется \term{квадратичным невычетом} (quadratic non-residue). Это различение имеет особое значение в наших исследованиях эллиптических кривых (\chaptname{} \ref{chap:elliptic_curves}), поскольку только квадратные числа могут фактически быть точками на эллиптической кривой.

Чтобы сделать интуицию о квадратичных вычетах и их корнях точной, мы даём следующее определение:

\begin{definition}
 Пусть $p \in \Prim $ -- простое число и $\F_p $ -- соответствующее простое поле. Тогда число $x\in \F_p$ называется \textbf{квадратным корнем} другого числа $y\in\F_p$, если $x$ является решением следующего уравнения:
\begin{equation}
x^2 = y
\end{equation}
В этом случае $y$ называется \term{квадратичным вычетом} (quadratic residue). С другой стороны, если дано $y$ и квадратное уравнение не имеет решения $x$, мы называем $ y $ \term{квадратичным невычетом} (quadratic non-residue).\footnote{Более подробное введение в квадратичные вычеты и их квадратные корни вместе с введением в алгоритмы, вычисляющие квадратные корни, можно найти, например, в \chaptname{} 1, \secname{} 1.5 книги \cite{cohen-2010}.}
\end{definition}

Для любого $ y \in \F_p $ мы обозначаем множество всех квадратных корней $ y $ в простом поле $ \F_p $ следующим образом:
\begin{equation}
\label{ded:square_root}
\sqrt{y}: = \{x \in \F_p \; | \; x^2 = y \}
\end{equation}
Неформально говоря, квадратичные вычеты -- это числа, имеющие квадратный корень, а квадратичные невычеты -- числа, не имеющие квадратных корней. Ситуация поэтому параллельна знакомому случаю целых чисел, где некоторые целые числа, например $4$ или $9$, имеют квадратный корень, а другие, например $2$ или $3$, не имеют (в кольце целых чисел).

Если $ y $ -- квадратичный невычет, то $ \sqrt{y} = \emptyset $ (пустое множество), а если $ y = 0 $, то $ \sqrt{y} = \{0 \} $.

Кроме того, если $y\neq 0$ является квадратичным вычетом, то оно имеет ровно два корня $\sqrt{y}=\{x,p-x\}$ для некоторого $x\in \F_p$. Мы принимаем соглашение называть меньший из них (при интерпретации как целого числа) \term{положительным квадратным корнем} (positive square root), а больший -- \term{отрицательным квадратным корнем} (negative square root).

Если $ p \in \Prim_{\geq 3} $ -- нечётное простое число с соответствующим простым полем $\F_p$, то существует ровно $(p+1)/2$ квадратичных вычетов и $(p-1)/2$ квадратичных невычетов.

\begin{example} [Квадратичные вычеты и корни в $ \F_5 $]
\label{example:quadratic_residue_F5}
Рассмотрим снова простое поле $\F_5$ из \examplename{} \ref{primfield_z_5}. Все квадратные числа можно найти на главной диагонали таблицы умножения в \eqref{Z5_tables}. Как видно, в $ \F_5 $ только числа $ 0 $, $ 1 $ и $ 4 $ имеют квадратные корни: $ \sqrt{0} = \{0 \} $, $ \sqrt{1} = \{1,4 \} $, $ \sqrt{2} = \emptyset $, $ \sqrt{3} = \emptyset $ и $ \sqrt{4} = \{2,3 \} $. Числа $0$, $1$ и $4$ поэтому являются квадратичными вычетами, а числа $2$ и $3$ -- квадратичными невычетами.
\end{example}
Чтобы описать, является ли элемент простого поля квадратным числом или нет, в литературе иногда встречается так называемый \term{символ Лежандра} (Legendre symbol) (например, \chaptname{} 1, \secname{} 1.5. книги \cite{cohen-2010}), определяемый следующим образом:

Пусть $ p \in \Prim $ -- простое число и $ y \in \F_p $ -- элемент из соответствующего простого поля. Тогда \textit{символ Лежандра} $ y $ определяется следующим образом:
\begin{equation}
\label{eq: Legendre-symbol}
\left (\frac{y}{p} \right): =
\begin{cases}
1 & \text{если $ y $ имеет квадратные корни} \\
-1 & \text{если $ y $ не имеет квадратных корней} \\
0 & \text{если $ y = 0 $}
\end{cases}
\end{equation}
\begin{example}
Рассмотрев снова квадратичные вычеты и невычеты в $\F_5$ из \examplename{} \ref{primfield_z_5}, мы можем вывести следующие символы Лежандра на основе \examplename{} \ref{example:quadratic_residue_F5}.

$$
\begin{array}{ccccc}
\left (\frac{0}{5} \right) = 0, &
\left (\frac{1}{5} \right) = 1, &
\left (\frac{2}{5} \right) = -1, &
\left (\frac{3}{5} \right) = -1, &
\left (\frac{4}{5} \right) = 1 \;.
\end{array}
$$
\end{example}
Символ Лежандра предоставляет критерий для решения, имеет ли элемент простого поля квадратичный корень или нет. Это, однако, не просто теоретическое применение: так называемый \term{критерий Эйлера} (Euler criterion) предоставляет компактный способ фактически вычислить символ Лежандра. Чтобы увидеть это, пусть $ p \in \Prim_{\geq 3} $ -- нечётное простое число и $ y \in \F_p $. Тогда символ Лежандра может быть вычислен следующим образом:
\begin{equation}
\label{eq: Euler_criterion}
\left (\frac{y}{p} \right) = y^{\frac{p-1}{2}}
\end{equation}
\begin{example}
Рассмотрев снова квадратичные вычеты и невычеты в $\F_5$ из \examplename{} \ref{example:quadratic_residue_F5}, мы можем вычислить следующие символы Лежандра с использованием критерия Эйлера:
\begin{align*}
\left (\frac{0}{5} \right) &= 0^{\frac{5-1}{2}}= 0^2=0\\
\left (\frac{1}{5} \right) &= 1^{\frac{5-1}{2}}= 1^2=1\\
\left (\frac{2}{5} \right) &= 2^{\frac{5-1}{2}}= 2^2=4 = -1\\
\left (\frac{3}{5} \right) &= 3^{\frac{5-1}{2}}= 3^2=4 =-1\\
\left (\frac{4}{5} \right) &= 4^{\frac{5-1}{2}}= 4^2=1
\end{align*}
\end{example}
\begin{exercise}
Рассмотрим простое поле $\F_{13}$ из \exercisename{} \ref{prime_field_F13}. Вычислите символ Лежандра $\left(\frac{x}{13} \right)$ и множество корней $\sqrt{x}$ для всех элементов $x\in \F_{13}$.
\end{exercise}
\subsubsection{Хеширование в простые поля}\label{hashing-prime-fields}
Важной задачей в криптографии является возможность хеширования в (различные подмножества) эллиптических кривых. Как мы увидим в \chaptname{} \ref{chap:elliptic_curves}, эти кривые часто определены над простыми полями, и хеширование в кривую может начинаться с хеширования в простое поле. Поэтому важно понимать, как хешировать в простые поля.

В \secname{} \ref{hash-to-modular-arithmetics}, мы рассмотрели несколько методов хеширования в кольца модульной арифметики $\Z_n$ для произвольного $n>1$. Поскольку простые поля -- это лишь частные случаи этих колец, все методы хеширования в $\Z_n$ могут использоваться и для хеширования в простые поля.

\subsection{Расширения простых полей}\label{field-extension}
Простые поля, как определено в предыдущем разделе, являются базовыми строительными блоками криптографии. Однако, как мы увидим в \chaptname{} \ref{chapter:zk-protocols}, так называемые SNARK-системы на основе паринга критически зависят от определённых групповых парингов \eqref{pairing-map}, определённых на эллиптических кривых над \term{расширениями простых полей} (prime field extensions). В этом разделе мы поэтому вводим эти расширения.\footnote{Более подробное введение можно найти, например, в \chaptname{} 2 книги \cite{nieder-1986}.}

Дано некоторое простое число $p\in \Prim$, натуральное число $m\in\N$ и неприводимый многочлен $P\in \F_p[x]$ степени $m$ с коэффициентами из простого поля $\F_p$, расширение простого поля $(\F_{p^m},+,\cdot)$ определяется следующим образом.

Множество $\F_{p^m}$ расширения простого поля задаётся множеством всех многочленов степени меньше $m$:
\begin{equation}
\label{eq:prime-extension-field}
\F_{p^m}:= \{a_{m-1} x^{m-1}+a_{m-2}x^{m-2}+\ldots+a_1 x+a_0\;|\; a_i\in \F_p\}
\end{equation}
Закон сложения расширения простого поля $\F_{p^m}$ задаётся обычным сложением многочленов, как определено в \eqref{def:polynomial_arithmetic}:
\begin{equation}
+:\; \F_{p^m}\times \F_{p^m} \to \F_{p^m}\; , \textstyle (\sum_{j=0}^m-1 a_j x^j,\sum_{j=0}^m-1 b_j x^j)\mapsto \sum_{j=0}^m-1 (a_j+b_j) x^j
\end{equation}
Закон умножения расширения простого поля $\F_{p^m}$ задаётся сначала умножением двух многочленов, как определено в \eqref{def:polynomial_arithmetic_mul}, а затем делением результата на неприводимый многочлен $P$ и сохранением остатка:
\begin{equation}
\cdot : \; \F_{p^m}\times \F_{p^m} \to \F_{p^m}\; , \textstyle\; (\sum_{j=0}^m a_j x^j,\sum_{j=0}^m b_j x^j)\mapsto \Zmod{\left(\sum _{n = 0} ^{2m} \sum _{i = 0} ^{n}{a} _{i }{{b} _{n-i}}{x} ^{n}\right)}{P}
\end{equation}
Нейтральным элементом аддитивной группы $(\F_{p^m},+)$ является нулевой многочлен $0$. Аддитивный обратный задаётся многочленом со всеми отрицательными коэффициентами. Нейтральным элементом мультипликативной группы $(\F^*_{p^m},\cdot)$ является единичный многочлен $1$. Мультипликативный обратный может быть вычислен расширенным алгоритмом Евклида (см. \ref{eq: erw_Eukl_algo})\sme{check reference}.

Мы видим из определения $\F_{p^m}$, что поле имеет характеристику $p$, поскольку мультипликативный нейтральный элемент $1$ эквивалентен мультипликативному элементу $1$ из лежащего в основе простого поля, и следовательно $\sum_{j=0}^p 1=0$. Кроме того, $\F_{p^m}$ конечно и содержит $p^m$ элементов, поскольку элементы являются многочленами степени $<m$, и каждый коэффициент $a_j$ может принимать $p$ различных значений. В дополнение, мы видим, что простое поле $\F_p$ является подполем $\F_{p^m}$, которое возникает, когда мы ограничиваем элементы $\F_{p^m}$ многочленами степени ноль.

Одним ключевым моментом является то, что построение $\F_{p^m}$ зависит от выбора неприводимого многочлена, и, по сути, различные выборы дадут различные таблицы умножения, поскольку остатки от деления произведения многочленов на эти многочлены будут различными.

Однако можно показать, что поля для различных выборов $P$ \term{изоморфны} (isomorphic), то есть между всеми ними существует взаимно однозначное соответствие. В результате, с абстрактной точки зрения, они являются одним и тем же. С точки зрения реализации, однако, некоторые выборы предпочтительнее других, поскольку они позволяют выполнять более быстрые вычисления.

\begin{remark}
Аналогично тому, как простые поля $\F_p$ генерируются начиная с кольца целых чисел, а затем делением на простое число $p$ и сохранением остатка, расширения простых полей $\F_{p^m}$ генерируются начиная с кольца $\F_p[x]$ многочленов, а затем делением их на неприводимый многочлен степени $m$ и сохранением остатка.

На самом деле, можно показать, что $\F_{p^m}$ -- это множество всех остатков при делении всех многочленов $Q\in \F_p[x]$ на неприводимый многочлен $P$ степени $m$. Это аналогично тому, как $\F_p$ -- это множество всех остатков при делении целых чисел на $p$.
\end{remark}

Любое поле $\F_{p^m}$, построенное описанным выше способом, является расширением поля $\F_p$. Более общо, поле $\F_{p^{m_2}}$ является расширением поля $\F_{p^{m_1}}$ тогда и только тогда, когда $m_1$ делит $m_2$. Из этого мы можем вывести, что для любого данного фиксированного простого числа существуют вложенные последовательности подполей всякий раз, когда степень $m_j$ делит степень $m_{j+1}$:

\begin{equation}
\F_p \subset \F_{p^{m_1}} \subset \cdots \subset \F_{p^{m_k}}
\end{equation}

Чтобы получить более интуитивную картину, мы построим расширение поля простого поля $\F_3$ в следующем примере, и мы увидим, как $\F_3$ располагается внутри этого расширения поля.
\begin{example}[Расширение поля $\F_{3^2}$]
\label{finite_field_F3_2}
В \exercisename{} \ref{exercise:F3} мы построили простое поле $\F_3$. В этом примере мы применяем определение расширения поля \eqref{eq:prime-extension-field} для построения расширения поля $\F_{3^2}$. Мы начинаем с выбора неприводимого многочлена степени $2$ с коэффициентами в $\F_3$. Пробуем
$P(t)=t^2+1$. Возможно, самый быстрый способ показать, что $P$ действительно неприводим, -- это просто подставить все элементы из $\F_3$, чтобы проверить, будет ли результат когда-либо равен нулю. Мы вычисляем следующим образом:
\begin{align*}
P(0) = 0^2+1 &= 1\\
P(1) = 1^2+1 &= 2\\
P(2) = 2^2+1 &=  1+1  = 2
\end{align*}
Это означает, что $P$ неприводим, поскольку все множители должны иметь форму $(t-a)$, где $a$ -- корень $P$. Множество $\F_{3^2}$ содержит все многочлены степени меньше двух с коэффициентами в $\F_3$, которые перечислены ниже:
\begin{equation}
\F_{3^2}=\{0,1,2,t,t+1,t+2,2t,2t+1,2t+2\}
\end{equation}

Как и ожидалось, наше расширение поля содержит $9$ элементов. Сложение определяется как сложение многочленов; например $(t+2) + (2t+2)= (1+2)t +(2+2)= 1$. Выполнение этого вычисления для всех элементов даёт следующую таблицу сложения
\begin{equation}\label{F32-add-table}
  \begin{tabular}{c | c c c c c c c c c}
    + & 0    & 1    & 2    & t    & t+1  & t+2  & 2t   & 2t+1 & 2t+2 \\\\\hline
    0 & 0    & 1    & 2    & t    & t+1  & t+2  & 2t   & 2t+1 & 2t+2 \\
    1 & 1    & 2    & 0    & t+1  & t+2  & t    & 2t+1 & 2t+2 & 2t   \\
    2 & 2    & 0    & 1    & t+2  & t    & t+1  & 2t+2 & 2t   & 2t+1 \\
    t & t    & t+1  & t+2  & 2t   & 2t+1 & 2t+2 & 0    & 1    & 2    \\
  t+1 & t+1  & t+2  & t    & 2t+1 & 2t+2 & 2t   & 1    & 2    & 0    \\
  t+2 & t+2  & t    & t+1  & 2t+2 & 2t   & 2t+1 & 2    & 0    & 1    \\
   2t & 2t   & 2t+1 & 2t+2 & 0    & 1    & 2    & t    & t+1  & t+2  \\
 2t+1 & 2t+1 & 2t+2 & 2t   & 1    & 2    & 0    & t+1  & t+2  & t    \\
 2t+2 & 2t+2 & 2t   & 2t+1 & 2    & 0    & 1    & t+2  & t    & t+1
  \end{tabular}
\end{equation}

Как мы видим, группа $(\F_3,+)$ является подгруппой группы $(\F_{3^2},+)$, полученной рассмотрением только первых трёх строк и столбцов этой таблицы.

Мы можем использовать таблицу сложения \eqref{F32-add-table}, чтобы вывести аддитивный обратный (отрицательный) любого элемента из $\F_{3^2}$. Например, в $\F_{3^2}$ мы имеем $-(2t+1)= t+2$, поскольку $(2t+1) + (t+2)=0$.

Умножение требует немного больше вычислений, поскольку мы сначала должны умножить многочлены, и всякий раз, когда результат имеет степень $\geq 2$, мы должны применить алгоритм деления многочленов (\algname{} \ref{alg_polynom_euclid_alg}), чтобы разделить произведение на многочлен $P$ и сохранить остаток. Чтобы увидеть, как это работает, вычислим произведение $t+2$ и $2t+2$ в $\F_{3^2}$:
\begin{align*}
(t+2) \cdot (2t+2) &= \Zmod{(2t^2 + 2t + t + 1)}{(t^2+1)} \\
                   &= \Zmod{(2t^2+1)}{(t^2+1)} & \#\; 2t^2+1:t^2+1= 2 + \frac{2}{t^2+1} \\
                   &= 2
\end{align*}
Это означает, что произведение $t+2$ и $2t+2$ в $\F_{3^2}$ равно $2$. Выполнение этого вычисления для всех элементов даёт следующую таблицу умножения:
\begin{equation}\label{F32-mult-table}
  \begin{tabular}{c | c c c c c c c c c}
$\cdot$ & 0    & 1    & 2    & t    & t+1  & t+2  & 2t   & 2t+1 & 2t+2 \\\\\hline
      0 & 0    & 0    & 0    & 0    & 0    & 0    & 0    & 0    & 0 \\
      1 & 0    & 1    & 2    & t    & t+1  & t+2  & 2t   & 2t+1 & 2t+2\\
      2 & 0    & 2    & 1    & 2t   & 2t+2 & 2t+1 & t    & t+2  & t+1 \\
      t & 0    & t    & 2t   & 2    & t+2  & 2t+2 & 1    & t+1  & 2t+1  \\
    t+1 & 0    & t+1  & 2t+2 & t+2  & 2t   & 1    & 2t+1 & 2    & t   \\
    t+2 & 0    & t+2  & 2t+1 & 2t+2 & 1    & t    & t+1  & 2t   & 2    \\
     2t & 0    & 2t   & t    & 1    & 2t+1 & t+1  & 2  & 2t+2 & t+2\\
   2t+1 & 0    & 2t+1 & t+2  & t+1  & 2    & 2t   & 2t+2 & t    & 1    \\
   2t+2 & 0    & 2t+2 & t+1  & 2t+1 & t    & 2    & t+2  & 1     & 2t
  \end{tabular}
\end{equation}
Как и в предыдущих примерах, мы можем использовать таблицу \eqref{F32-mult-table}, чтобы вывести мультипликативный обратный любого ненулевого элемента из $\F_{3^2}$. Например, в $\F_{3^2}$ мы имеем $(2t+1)^{-1}= 2t+2 $, поскольку $(2t+1) \cdot (2t+2)=1$.

Глядя на таблицу умножения \eqref{F32-mult-table}, мы также можем видеть, что единственными квадратичными вычетами в $\F_{3^2}$ являются элементы из множества $\{0,1,2, t, 2t\}$, с
$\sqrt{0}=\{0\}$, $\sqrt{1}=\{1,2\}$, $\sqrt{2}=\{t, 2t\}$, $\sqrt{t}=\{t+2,2t+1\}$ и $\sqrt{2t}=\{t+1,2t+2\}$.

Поскольку $\F_{3^2}$ -- поле, мы можем решать уравнения, как мы делали бы для других полей (таких как рациональные числа). Чтобы увидеть это, найдём все $x\in\F_{3^2}$, которые решают квадратное уравнение $(t+1)(x^2 + (2t+2)) = 2$. Мы вычисляем следующим образом:
\begin{align*}
(t+1)(x^2 + (2t+2))    &= 2 &\text{\# 2-й дистрибутивный закон}\\
(t+1)x^2 + (t+1)(2t+2) &= 2 \\
(t+1)x^2 + (t)         &= 2 &\text{\# 2 прибавляем аддитивный обратный $t$}\\
(t+1)x^2 + (t) + (2t)  &= (2) + (2t) \\
(t+1)x^2               &= 2t+2 & \text{\# умножаем на мультипликативный обратный $t+1$}\\
(t+2)(t+1)x^2          &=(t+2)(2t+2) & \text{\# умножаем на мультипликативный обратный $t+1$}\\
x^2                    &= 2 & \text{\# 2 -- квадратичный вычет. Берём корни.}\\
x &\in \{t, 2t\}
\end{align*}
Вычисления в расширениях полей, вероятно, находятся на грани того, что можно разумно делать бумагой и ручкой. К счастью, Sage предоставляет нам простой способ выполнять эти вычисления.
\begin{sagecommandline}
sage: Z3 = GF(3) # простое поле
sage: Z3t.<t> = Z3[] # многочлены над Z3
sage: P = Z3t(t^2+1)
sage: P.is_irreducible()
sage: F3_2.<t> = GF(3^2, name='t', modulus=P) # Расширение поля F_3^2
sage: F3_2
sage: F3_2(t+2)*F3_2(2*t+2) == F3_2(2)
sage: F3_2(2*t+2)^(-1) # мультипликативный обратный
sage: # проверяем наше решение (t+1)(x^2 + (2t+2)) = 2
sage: F3_2(t+1)*(F3_2(t)**2 + F3_2(2*t+2)) == F3_2(2)
sage: F3_2(t+1)*(F3_2(2*t)**2 + F3_2(2*t+2)) == F3_2(2)
\end{sagecommandline}
\end{example}
\begin{exercise}
Рассмотрим расширение поля $\F_{3^2}$ из предыдущего примера и найдите все пары элементов $(x,y)\in\F_{3^2}$, для которых выполняется следующее уравнение:

\begin{equation}
y^2 = x^3 + 4
\end{equation}
\end{exercise}

\begin{exercise} Покажите, что многочлен $Q=x^2+x+2$ из $\F_3[x]$ неприводим. Постройте таблицу умножения $\F_{3^2}$ относительно $Q$ и сравните её с таблицей умножения $\F_{3^2}$ из примера \ref{finite_field_F3_2}.
\end{exercise}

\begin{exercise} Покажите, что многочлен $P=t^3+t+1$ из $\F_5[t]$ неприводим. Затем рассмотрите расширение поля $\F_{5^3}$, определённое относительно $P$. Вычислите мультипликативный обратный $(2t^2+4)\in\F_{5^3}$ с использованием расширенного алгоритма Евклида. Затем найдите все $x\in\F_{5^3}$, которые решают следующее уравнение:
\begin{equation}
(2t^2+4)(x-(t^2+4t+2))= (2t+3)
\end{equation}
\end{exercise}

\begin{exercise}
\label{exercise:finite_fieldF5_2} Рассмотрим простое поле $\F_5$. Покажите, что многочлен $P=x^2+2$ из $\F_5[x]$ неприводим. Реализуйте конечное поле $\F_{5^2}$ в Sage.
\end{exercise}

\begin{comment}
\subsection{Хеширование в поля расширения} On page \pageref{hashing-prime-fields},\sme{check reference} we have seen how to hash into prime fields. As elements of extension fields can be seen as polynomials over prime fields, hashing into extension fields is therefore possible if every coefficient of the polynomial is hashed independently.
\end{comment}
\section{Проективные плоскости}\label{sec:planes}
Проективные плоскости -- это особые геометрические конструкции, определённые над данным полем. В некотором смысле, проективные плоскости расширяют концепцию обычной евклидовой плоскости, включая <<точки на бесконечности>>.\footnote{Подробное объяснение идей, приводящих к определению проективных плоскостей, можно найти, например, в \chaptname{} 2 книги \cite{ellis-1992} или в приложении A книги \cite{silverman-1994}.}

Чтобы понять идею построения проективных плоскостей, отметим, что на обычной евклидовой плоскости две прямые либо пересекаются в одной точке, либо параллельны. В последнем случае обе прямые либо совпадают, то есть пересекаются во всех точках, либо не пересекаются вообще. Проективную плоскость можно тогда представить как обычную плоскость, но снабжённую дополнительной <<точкой на бесконечности>> так, что две различные прямые всегда пересекаются в одной точке. Параллельные прямые пересекаются <<на бесконечности>>.

Такое включение точек на бесконечности делает проективные плоскости особенно полезными при описании эллиптических кривых, поскольку описание такой кривой на обычной плоскости требует дополнительного символа для <<точки на бесконечности>>, чтобы дать множеству точек на кривой структуру группы \ref{sec:short_weierstrass_curve}. Перевод кривой в проективную геометрию включает эту <<точку на бесконечности>> более естественным образом в множество всех точек на проективной плоскости.

Более точно, пусть $\F$ -- поле, $\F^3:=\F\times \F\times \F$ -- множество всех кортежей из трёх элементов над $\F$ и $x\in \F^3$ с $x=(X,Y,Z)$. Тогда существует ровно одна \textit{прямая} $L_x$ в $\F^3$, которая пересекает и $(0,0,0)$, и $x$, задаваемая множеством $L_x=\{(k\cdot X,k\cdot Y, k\cdot Z)\;|\; k\in\F\}$. Точка на \textbf{проективной плоскости} над $\F$ может быть определена как такая \term{прямая} (line), если исключить пересечение этой прямой с $(0,0,0)$. Это приводит к следующему определению \term{точки} (point) в проективной геометрии:
\begin{equation}
\label{def:projective_coordinate}
[X:Y:Z] := \{(k\cdot X,k\cdot Y, k\cdot Z)\;|\; k\in\F^*\}
\end{equation}
Точки в проективной геометрии поэтому являются прямыми в $\F^3$, где пересечение с $(0,0,0)$ исключено. Данное поле $\F$, \term{проективная плоскость} (projective plane) этого поля определяется как множество всех точек, исключая точку $[0:0:0]$:
\begin{equation}
\F\mathbb{P}^2:=\{[X:Y:Z]\;|\; (X,Y,Z)\in \F^3\text{ при } (X,Y,Z)\neq (0,0,0)\}
\end{equation}
Можно показать, что проективная плоскость над конечным полем $\F_{p^m}$ содержит $p^{2m}+p^m+1$ элементов.

Чтобы понять, почему проективная точка $[X:Y:Z]$ также является прямой, рассмотрим ситуацию, когда лежащее в основе поле $\F$ является множеством рациональных чисел $\Q$. В этом случае $\Q^3$ может быть рассмотрено как трёхмерное пространство, а $[X:Y:Z]$ -- обычная прямая в этом трёхмерном пространстве, которая пересекает нуль и точку с координатами $X$, $Y$ и $Z$ так, что пересечение с нулём исключено.

Ключевое наблюдение здесь заключается в том, что точки на проективной плоскости $\F\mathbb{P}^2$ являются прямыми в трёхмерном пространстве $\F^3$. Однако следует помнить, что для конечных полей термины \term{пространство} (space) и \term{прямая} (line) имеют очень мало визуального сходства со своими аналогами над множеством рациональных чисел.

Из этого следует, что точки $[X:Y:Z]\in \F\mathbb{P}^2$ описываются не просто фиксированными координатами $(X,Y,Z)$, а \term{множествами координат} (sets of coordinates), где две различные координаты $(X_1,Y_1,Z_1)$ и $(X_2,Y_2,Z_2)$ описывают одну и ту же точку тогда и только тогда, когда существует некоторый ненулевой элемент поля $k\in\F^*$ такой, что $(X_1,Y_1,Z_1) = (k\cdot X_2,k\cdot Y_2,k\cdot Z_2)$. Точки $[X:Y:Z]$ называются \term{проективными координатами} (projective coordinates).

\begin{notation}[Проективные координаты]
Проективные координаты формы $[X:Y:1]$ являются описаниями так называемых \textbf{аффинных точек} (affine points). Проективные координаты формы $[X:Y:0]$ являются описаниями так называемых \textbf{точек на бесконечности} (points at infinity). В частности, проективная координата $[1:0:0]$ описывает так называемую \textbf{прямую на бесконечности} (line at infinity).
\end{notation}
\begin{example} Рассмотрим поле $\F_3$ из \exercisename{} \ref{exercise:F3} . Поскольку это поле содержит только три элемента, построение соответствующей проективной плоскости $\F_3\mathbb{P}^2$ не требует слишком больших усилий, и мы знаем, что она содержит только $13$ элементов.

Чтобы найти $\F_3\mathbb{P}^2$, мы должны вычислить множество всех прямых в $\F_3\times \F_3\times \F_3$, пересекающих $(0,0,0)$, исключая их пересечение с $(0,0,0)$. Поскольку эти прямые параметризуются кортежами $(x_1,x_2,x_3)$, мы вычисляем следующим образом:
\begin{align*}
[0:0:1] &= \{(k\cdot x_1,k\cdot x_2, k\cdot x_3)\;|\; k\in\F_3^*\}
          = \{(0,0,1), (0,0,2)\}\\
[0:0:2] &= \{(k\cdot x_1,k\cdot x_2, k\cdot x_3)\;|\; k\in\F_3^*\}
          = \{(0,0,2), (0,0,1)\}
          = [0:0:1]\\
[0:1:0] &= \{(k\cdot x_1,k\cdot x_2, k\cdot x_3)\;|\; k\in\F_3^*\}
          = \{(0,1,0), (0,2,0)\}\\
[0:1:1] &= \{(k\cdot x_1,k\cdot x_2, k\cdot x_3)\;|\; k\in\F_3^*\}
          = \{(0,1,1), (0,2,2)\}\\
[0:1:2] &= \{(k\cdot x_1,k\cdot x_2, k\cdot x_3)\;|\; k\in\F_3^*\}
          = \{(0,1,2), (0,2,1)\}\\
[0:2:0] &= \{(k\cdot x_1,k\cdot x_2, k\cdot x_3)\;|\; k\in\F_3^*\}
          = \{(0,2,0), (0,1,0)\}
          = [0:1:0]\\
[0:2:1] &= \{(k\cdot x_1,k\cdot x_2, k\cdot x_3)\;|\; k\in\F_3^*\}
          = \{(0,2,1), (0,1,2)\}
          = [0:1:2]\\
[0:2:2] &= \{(k\cdot x_1,k\cdot x_2, k\cdot x_3)\;|\; k\in\F_3^*\}
          = \{(0,2,2), (0,1,1)\}
          = [0:1:1]\\
[1:0:0] &= \{(k\cdot x_1,k\cdot x_2, k\cdot x_3)\;|\; k\in\F_3^*\}
          = \{(1,0,0), (2,0,0)\}\\
[1:0:1] &= \{(k\cdot x_1,k\cdot x_2, k\cdot x_3)\;|\; k\in\F_3^*\}
          = \{(1,0,1), (2,0,2)\}\\
[1:0:2] &= \{(k\cdot x_1,k\cdot x_2, k\cdot x_3)\;|\; k\in\F_3^*\}
          = \{(1,0,2), (2,0,1)\}\\
[1:1:0] &= \{(k\cdot x_1,k\cdot x_2, k\cdot x_3)\;|\; k\in\F_3^*\}
          = \{(1,1,0), (2,2,0)\}\\
[1:1:1] &= \{(k\cdot x_1,k\cdot x_2, k\cdot x_3)\;|\; k\in\F_3^*\}
          = \{(1,1,1), (2,2,2)\}\\
[1:1:2]&= \{(k\cdot x_1,k\cdot x_2, k\cdot x_3)\;|\; k\in\F_3^*\}
          = \{(1,1,2), (2,2,1)\}\\
[1:2:0] &= \{(k\cdot x_1,k\cdot x_2, k\cdot x_3)\;|\; k\in\F_3^*\}
          = \{(1,2,0), (2,1,0)\}\\
[1:2:1] &= \{(k\cdot x_1,k\cdot x_2, k\cdot x_3)\;|\; k\in\F_3^*\}
          = \{(1,2,1), (2,1,2)\}\\
[1:2:2] &= \{(k\cdot x_1,k\cdot x_2, k\cdot x_3)\;|\; k\in\F_3^*\}
          = \{(1,2,2), (2,1,1)\}\\
[2:0:0] &= \{(k\cdot x_1,k\cdot x_2, k\cdot x_3)\;|\; k\in\F_3^*\}
          = \{(2,0,0), (1,0,0)\}
          = [1:0:0]\\
[2:0:1] &= \{(k\cdot x_1,k\cdot x_2, k\cdot x_3)\;|\; k\in\F_3^*\}
          = \{(2,0,1), (1,0,2)\}
          = [1:0:2]\\
[2:0:2] &= \{(k\cdot x_1,k\cdot x_2, k\cdot x_3)\;|\; k\in\F_3^*\}
          = \{(2,0,2), (1,0,1)\}
          = [1:0:1]\\
[2:1:0] &= \{(k\cdot x_1,k\cdot x_2, k\cdot x_3)\;|\; k\in\F_3^*\}
          = \{(2,1,0), (1,2,0)\}
          = [1:2:0]\\
[2:1:1] &= \{(k\cdot x_1,k\cdot x_2, k\cdot x_3)\;|\; k\in\F_3^*\}
          = \{(2,1,1), (1,2,2)\}
          = [1:2:2]\\
[2:1:2] &= \{(k\cdot x_1,k\cdot x_2, k\cdot x_3)\;|\; k\in\F_3^*\}
          = \{(2,1,2), (1,2,1)\}
          = [1:2:1]\\
[2:2:0] &= \{(k\cdot x_1,k\cdot x_2, k\cdot x_3)\;|\; k\in\F_3^*\}
          = \{(2,2,0), (1,1,0)\}
          = [1:1:0]\\
[2:2:1] &= \{(k\cdot x_1,k\cdot x_2, k\cdot x_3)\;|\; k\in\F_3^*\}
          = \{(2,2,1), (1,1,2)\}
          = [1:1:2]\\
[2:2:2] &= \{(k\cdot x_1,k\cdot x_2, k\cdot x_3)\;|\; k\in\F_3^*\}
          = \{(2,2,2), (1,1,1)\}
          = [1:1:1]
\end{align*}
Эти прямые определяют $13$ точек в проективной плоскости $\F_3\mathbb{P}$:
\begin{multline*}
\F_3\mathbb{P} = \{ [0:0:1], [0:1:0], [0:1:1], [0:1:2], [1:0:0], [1:0:1], \\ [1:0:2], [1:1:0], [1:1:1], [1:1:2], [1:2:0], [1:2:1], [1:2:2]\}
\end{multline*}
Эта проективная плоскость содержит $9$ аффинных точек, три точки на бесконечности и одну прямую на бесконечности.

Чтобы лучше понять неоднозначность проективных координат, рассмотрим точку $[1:2:2]$. Поскольку эта точка в проективной плоскости является прямой в $\F_3^3\backslash\{(0,0,0)\}$, она имеет проективные координаты $(1,2,2)$, а также $(2,1,1)$, поскольку первые координаты дают вторые при умножении в $\F_3$ на множитель $2$. Кроме того, отметим, что по тем же причинам точки $[1:2:2]$ и $[2:1:1]$ совпадают, так как их лежащие в основе множества равны.
\end{example}
\begin{exercise}
Постройте так называемую \term{плоскость Фано} (Fano plane), то есть проективную плоскость над конечным полем $\F_2$.
\end{exercise}
